\documentclass[12pt,a4paper]{article}

\usepackage[a4paper,margin=1in]{geometry}
\usepackage{amsmath,amssymb,bm,mathtools}
\usepackage{physics}
\usepackage{siunitx}
\usepackage{setspace}
\usepackage{microtype}
\usepackage{titlesec}
\usepackage{float}
\usepackage{graphicx}
\setstretch{1.15}
\setlength{\parindent}{0pt}
\setlength{\parskip}{6pt}

\titleformat{\section}{\large\bfseries}{\thesection}{0.7em}{}
\titleformat{\subsection}{\normalsize\bfseries}{\thesubsection}{0.7em}{}

\newcommand{\ketg}{\left|g\right\rangle}
\newcommand{\kete}{\left|e\right\rangle}
\newcommand{\Efield}{\mathbf E}
\newcommand{\rvec}{\mathbf r}

\begin{document}

\begin{center}
    {\Large\bfseries Channel Estimation for a Rydberg Atomic Receiver}\\[4pt]
\end{center}

\hrule
\vspace{8pt}

\textbf{Notation.} In the Schr\"odinger equation, $i\hbar$ uses
$i=j=\sqrt{-1}$ (the imaginary unit is written $j$ elsewhere). The
subscript $i$ always denotes the vapor-cell index.

% ============================================================
\section{Quantum System and Two-Level Rydberg State}
% ============================================================

A quantum system is described by a state vector. For the two-level
Rydberg-atom model, the ground state and excited state are denoted
by
\[
\ketg,
\qquad
\kete.
\]

The state of atom \(i\) is written as
\begin{equation}
\ket{\Phi_i(t)}
=
\zeta_{g,i}(t)\ketg
+
\zeta_{e,i}(t)\kete .
\tag{1}
\end{equation}

The quantities
\(\zeta_{g,i}(t)\) and \(\zeta_{e,i}(t)\) are the probability
amplitudes corresponding to the ground and excited states,
respectively.

The probability of finding the atom in the ground state is
\[
P_{g,i}(t)
=
\left|\zeta_{g,i}(t)\right|^2,
\]
while the probability of finding the atom in the excited state is
\[
P_{e,i}(t)
=
\left|\zeta_{e,i}(t)\right|^2.
\]

Since the atom must be in either the ground state or the excited
state, the total probability is unity:
\[
P_{g,i}(t)+P_{e,i}(t)=1.
\]

Therefore,
\[
\boxed{
\left|\zeta_{g,i}(t)\right|^2
+
\left|\zeta_{e,i}(t)\right|^2
=1.
}
\]

The two-level basis is chosen as
\[
\kete
=
\begin{bmatrix}
1\\
0
\end{bmatrix},
\qquad
\ketg
=
\begin{bmatrix}
0\\
1
\end{bmatrix}.
\]

Substituting these basis vectors into (1),
\[
\begin{aligned}
\ket{\Phi_i(t)}
&=
\zeta_{g,i}(t)
\begin{bmatrix}
0\\
1
\end{bmatrix}
+
\zeta_{e,i}(t)
\begin{bmatrix}
1\\
0
\end{bmatrix}\\
&=
\begin{bmatrix}
\zeta_{e,i}(t)\\
\zeta_{g,i}(t)
\end{bmatrix}.
\end{aligned}
\]

Hence, for the basis ordering
\[
(\kete,\ketg),
\]
the state vector is
\[
\ket{\Phi_i(t)}
=
\begin{bmatrix}
\zeta_{e,i}(t)\\
\zeta_{g,i}(t)
\end{bmatrix}.
\]

% ============================================================
\section{Time-Dependent Schr\"odinger Equation}
% ============================================================

The time evolution of a quantum state is described by the
time-dependent Schr\"odinger equation.

In the absence of an external interaction,
\[
i\hbar
\frac{\partial}{\partial t}
\ket{\Phi_i(t)}
=
\hat H_i
\ket{\Phi_i(t)},
\]
where \(\hat H_i\) is the free Hamiltonian of the atom.

When an external electromagnetic field interacts with the atom,
an interaction Hamiltonian \(\hat V_i\) is introduced. Therefore,
the total Hamiltonian is
\[
\hat H_{i,\mathrm{tot}}
=
\hat H_i+\hat V_i.
\]

Hence,
\begin{equation}
i\hbar
\frac{\partial}{\partial t}
\ket{\Phi_i(t)}
=
(\hat H_i+\hat V_i)
\ket{\Phi_i(t)}.
\tag{2}
\end{equation}

The free Hamiltonian describes the internal energy of the atom,
whereas \(\hat V_i\) represents the interaction between the atom
and the incident electromagnetic field.

% ============================================================
\section{Free Hamiltonian of the Two-Level Atom}
% ============================================================

The atomic energy states satisfy the energy eigenvalue equation
\[
\hat H_i\ket{\psi_n}
=
E_n\ket{\psi_n}.
\]

For the two states,
\[
\hat H_i\kete
=
E_e\kete,
\]
and
\[
\hat H_i\ketg
=
E_g\ketg.
\]

The energies of the two states are related to their angular
frequencies by
\[
E_e=\hbar\omega_e,
\qquad
E_g=\hbar\omega_g.
\]

Using the basis
\[
\kete=
\begin{bmatrix}
1\\
0
\end{bmatrix},
\qquad
\ketg=
\begin{bmatrix}
0\\
1
\end{bmatrix},
\]
write the free Hamiltonian in the general form
\[
\hat H_i
=
\begin{bmatrix}
a&b\\
c&d
\end{bmatrix}.
\]

Applying \(\hat H_i\) to the excited state gives
\[
\begin{aligned}
\hat H_i\kete
&=
\begin{bmatrix}
a&b\\
c&d
\end{bmatrix}
\begin{bmatrix}
1\\
0
\end{bmatrix}\\
&=
\begin{bmatrix}
a\\
c
\end{bmatrix}.
\end{aligned}
\]

From the energy eigenvalue equation,
\[
\hat H_i\kete
=
E_e\kete.
\]

Therefore,
\[
E_e\kete
=
E_e
\begin{bmatrix}
1\\
0
\end{bmatrix}
=
\begin{bmatrix}
E_e\\
0
\end{bmatrix}.
\]

Comparing the two expressions,
\[
\begin{bmatrix}
a\\
c
\end{bmatrix}
=
\begin{bmatrix}
E_e\\
0
\end{bmatrix},
\]
which gives
\[
a=E_e,
\qquad
c=0.
\]

Similarly, applying \(\hat H_i\) to the ground state,
\[
\begin{aligned}
\hat H_i\ketg
&=
\begin{bmatrix}
a&b\\
c&d
\end{bmatrix}
\begin{bmatrix}
0\\
1
\end{bmatrix}\\
&=
\begin{bmatrix}
b\\
d
\end{bmatrix}.
\end{aligned}
\]

The energy eigenvalue equation gives
\[
\hat H_i\ketg
=
E_g\ketg
=
E_g
\begin{bmatrix}
0\\
1
\end{bmatrix}
=
\begin{bmatrix}
0\\
E_g
\end{bmatrix}.
\]

Therefore,
\[
\begin{bmatrix}
b\\
d
\end{bmatrix}
=
\begin{bmatrix}
0\\
E_g
\end{bmatrix},
\]
and hence
\[
b=0,
\qquad
d=E_g.
\]

Thus,
\begin{equation}
\boxed{
\hat H_i
=
\begin{bmatrix}
E_e&0\\
0&E_g
\end{bmatrix}.
}
\tag{3}
\end{equation}

Using
\[
E_e=\hbar\omega_e,
\qquad
E_g=\hbar\omega_g,
\]
we obtain
\begin{equation}
\boxed{
\hat H_i
=
\begin{bmatrix}
\hbar\omega_e&0\\
0&\hbar\omega_g
\end{bmatrix}
=
\operatorname{diag}
(\hbar\omega_e,\hbar\omega_g).
}
\tag{4}
\end{equation}

% ============================================================
\section{One-Dimensional Vapor-Cell Array Geometry}
% ============================================================

Consider a one-dimensional array of vapor cells arranged along
the \(x\)-axis.

\begin{figure}[H]
    \centering
    \IfFileExists{1D_array_geometry.png}{%
    \includegraphics[width=0.85\textwidth]{1D_array_geometry.png}}{%
    \fbox{\parbox{0.8\textwidth}{\centering\vspace{2cm}
    [Place \texttt{1D\_array\_geometry.png} in this folder]\vspace{2cm}}}}
    \caption{1D uniform linear array geometry showing the inter-cell spacing $d$ and the position of the $i$-th vapor cell.}
    \label{fig:1D_array_geometry}
\end{figure}

Let the spacing between two adjacent vapor cells be \(d\), and
choose vapor cell \(1\) as the reference cell.

The position of vapor cell \(i\) is
\[
x_i=(i-1)d.
\]

Therefore,
\[
x_1=0,
\qquad
x_2=d,
\qquad
x_3=2d,
\qquad
\ldots
\]

The separation between the reference cell and cell \(i\) is
\[
x_i-x_1
=
(i-1)d.
\]

Consider user \(k\) and propagation path \(l\). Let
\(\theta_{l,k}\) denote the angle of arrival of this path measured
with respect to the array axis.

Under the far-field plane-wave assumption, the propagation paths
arriving at the different vapor cells are approximately parallel.

Hence,
\begin{equation}
\boxed{
\Delta r_{i,k,l}
=
(i-1)d\cos\theta_{l,k}.
}
\tag{5}
\end{equation}

where
\[
\Delta r_{i,k,l}
=
r_{i,k,l}-r_{1,k,l}.
\]

Here \(r_{i,k,l}\) denotes the propagation distance from user \(k\),
through path \(l\), to vapor cell \(i\).

For the second vapor cell,
\[
\Delta r_{2,k,l}
=
d\cos\theta_{l,k}.
\]

For the third vapor cell,
\[
\Delta r_{3,k,l}
=
2d\cos\theta_{l,k}.
\]

In general,
\[
\Delta r_{i,k,l}
=
(i-1)d\cos\theta_{l,k}.
\]

% ============================================================
\section{Propagation Phase Difference}
% ============================================================

Consider a sinusoidal plane wave represented by
\begin{equation}
E(r,t)
=
E_0
\cos
\left(
\omega t+kr+\phi_0
\right),
\tag{6}
\end{equation}
where
\[
k=\frac{2\pi}{\lambda}
\]
is the wavenumber.

If the propagation distance changes from \(r\) to
\(r+\Delta r\), the phase changes from
\[
kr
\]
to
\[
k(r+\Delta r).
\]

Therefore,
\[
\begin{aligned}
\Delta\phi
&=
k(r+\Delta r)-kr\\
&=
kr+k\Delta r-kr\\
&=
k\Delta r.
\end{aligned}
\]

Using
\[
k=\frac{2\pi}{\lambda},
\]
we obtain
\begin{equation}
\boxed{
\Delta\phi
=
\frac{2\pi}{\lambda}\Delta r.
}
\tag{7}
\end{equation}

For user \(k\), path \(l\), and vapor cell \(i\),
\[
\Delta r
=
\Delta r_{i,k,l}.
\]

Therefore,
\[
\begin{aligned}
\Delta\phi_{i,k,l}
&=
\frac{2\pi}{\lambda}
\Delta r_{i,k,l}\\
&=
\frac{2\pi}{\lambda}
(i-1)d\cos\theta_{l,k}\\
&=
(i-1)
\frac{2\pi d}{\lambda}
\cos\theta_{l,k}.
\end{aligned}
\]

Define the phase shift corresponding to one array spacing as
\begin{equation}
\boxed{
\phi_{l,k}
=
\frac{2\pi d}{\lambda}
\cos\theta_{l,k}.
}
\tag{8}
\end{equation}

Consequently,
\begin{equation}
\boxed{
\Delta\phi_{i,k,l}
=
(i-1)\phi_{l,k}.
}
\tag{9}
\end{equation}

For the reference vapor cell,
\[
i=1,
\]
and hence
\[
\Delta\phi_{1,k,l}=0.
\]

For the second vapor cell,
\[
i=2,
\]
and hence
\[
\Delta\phi_{2,k,l}=\phi_{l,k}.
\]

For the third vapor cell,
\[
i=3,
\]
and hence
\[
\Delta\phi_{3,k,l}=2\phi_{l,k}.
\]

Thus, the phase shift increases by \(\phi_{l,k}\) for every
additional array spacing.

% ============================================================
\section{Electric Field of One Propagation Path}
% ============================================================

Consider user \(k\) and propagation path \(l\).

The electric field corresponding to this propagation path can be
written as
\begin{equation}
\mathbf E_{k,l}(r,t)
=
\mathbf E_{0,k,l}
\cos
\left(
\omega t+kr+\phi_{0,k,l}
\right).
\tag{10}
\end{equation}

At vapor cell \(i\), the propagation distance is
\(r_{i,k,l}\). Therefore,
\begin{equation}
\mathbf E_{i,k,l}(t)
=
\mathbf E_{0,i,k,l}
\cos
\left(
\omega t
+
kr_{i,k,l}
+
\phi_{0,k,l}
\right).
\tag{11}
\end{equation}

Choose vapor cell \(1\) as the phase reference and define
\[
\Delta r_{i,k,l}
=
r_{i,k,l}-r_{1,k,l}.
\]

Therefore,
\[
r_{i,k,l}
=
r_{1,k,l}
+
\Delta r_{i,k,l}.
\]

Substituting this into (11),
\[
\begin{aligned}
\mathbf E_{i,k,l}(t)
&=
\mathbf E_{0,i,k,l}
\cos
\left[
\omega t
+
k
\left(
r_{1,k,l}+\Delta r_{i,k,l}
\right)
+
\phi_{0,k,l}
\right]\\
&=
\mathbf E_{0,i,k,l}
\cos
\left[
\omega t
+
k\Delta r_{i,k,l}
+
kr_{1,k,l}
+
\phi_{0,k,l}
\right].
\end{aligned}
\]

Define the common reference phase as
\[
\phi_{k,l}^{\mathrm{ref}}
=
kr_{1,k,l}
+
\phi_{0,k,l}.
\]

Hence,
\[
\mathbf E_{i,k,l}(t)
=
\mathbf E_{0,i,k,l}
\cos
\left[
\omega t
+
k\Delta r_{i,k,l}
+
\phi_{k,l}^{\mathrm{ref}}
\right].
\]

Taking vapor cell \(1\) as the phase reference gives
\[
\phi_{k,l}^{\mathrm{ref}}=0.
\]

Therefore,
\begin{equation}
\mathbf E_{i,k,l}(t)
=
\mathbf E_{0,i,k,l}
\cos
\left(
\omega t+k\Delta r_{i,k,l}
\right).
\tag{12}
\end{equation}

Using
\[
\Delta r_{i,k,l}
=
(i-1)d\cos\theta_{l,k},
\]
we obtain
\[
\mathbf E_{i,k,l}(t)
=
\mathbf E_{0,i,k,l}
\cos
\left[
\omega t
+
k(i-1)d\cos\theta_{l,k}
\right].
\]

Using
\[
k=\frac{2\pi}{\lambda},
\]
the electric field becomes
\[
\mathbf E_{i,k,l}(t)
=
\mathbf E_{0,i,k,l}
\cos
\left[
\omega t
+
(i-1)
\frac{2\pi d}{\lambda}
\cos\theta_{l,k}
\right].
\]

From (8),
\[
\phi_{l,k}
=
\frac{2\pi d}{\lambda}
\cos\theta_{l,k}.
\]

Therefore,
\begin{equation}
\boxed{
\mathbf E_{i,k,l}(t)
=
\mathbf E_{0,i,k,l}
\cos
\left[
\omega t+(i-1)\phi_{l,k}
\right].
}
\tag{13}
\end{equation}

% ============================================================
\section{Vector Representation and Polarization}
% ============================================================

The electric field is a vector quantity. Its amplitude can therefore
be separated into a scalar amplitude and a polarization direction.

Let
\[
\boldsymbol{\epsilon}_{i,k,l}
\]
denote the polarization vector corresponding to user \(k\),
path \(l\), and vapor cell \(i\).

The vector amplitude can be written as
\[
\mathbf E_{0,i,k,l}
=
E_{0,i,k,l}
\boldsymbol{\epsilon}_{i,k,l}.
\]

Substituting this into (13),
\[
\mathbf E_{i,k,l}(t)
=
E_{0,i,k,l}
\boldsymbol{\epsilon}_{i,k,l}
\cos
\left[
\omega t+(i-1)\phi_{l,k}
\right].
\]

Hence,
\begin{equation}
\boxed{
\mathbf E_{i,k,l}(t)
=
E_{0,i,k,l}
\boldsymbol{\epsilon}_{i,k,l}
\cos
\left[
\omega t+(i-1)\phi_{l,k}
\right].
}
\tag{14}
\end{equation}

The polarization vector represents the direction of the electric
field, while \(E_{0,i,k,l}\) represents its scalar amplitude.

% ============================================================
\section{Transmitted Symbol and Propagation Coefficient}
% ============================================================

Consider user \(k\) transmitting the symbol
\[
s_k.
\]

Let
\[
\alpha_{l,k}
\]
denote the complex coefficient associated with propagation path
\(l\) from user \(k\).

The scalar field amplitude is represented as
\[
E_{0,i,k,l}
=
\alpha_{l,k}s_k.
\]

Substituting this relation into (14),
\[
\mathbf E_{i,k,l}(t)
=
\alpha_{l,k}s_k
\boldsymbol{\epsilon}_{i,k,l}
\cos
\left[
\omega t+(i-1)\phi_{l,k}
\right].
\]

Therefore,
\begin{equation}
\boxed{
\mathbf E_{i,k,l}(t)
=
\boldsymbol{\epsilon}_{i,k,l}
\alpha_{l,k}s_k
\cos
\left[
\omega t+(i-1)\phi_{l,k}
\right].
}
\tag{15}
\end{equation}

The three terms in this expression represent the polarization
direction, propagation-path coefficient, and transmitted symbol,
respectively:
\[
\boldsymbol{\epsilon}_{i,k,l},
\qquad
\alpha_{l,k},
\qquad
s_k.
\]

The phase accumulated at vapor cell \(i\) is represented by
\[
(i-1)\phi_{l,k}.
\]

\textbf{Remark.} Equation (15) follows the form of Eq.~(3) in
Xu \emph{et al.}, where the complex baseband amplitude
$\alpha_{l,k}s_k$ multiplies the real carrier. A strictly real field is
$\mathbf E_{i,k,l}(t)=\operatorname{Re}\{\boldsymbol\epsilon_{i,k,l}\alpha_{l,k}s_k
e^{j[\omega t+(i-1)\phi_{l,k}]}\}$; using this form conjugates the
coupling coefficient obtained later in (40), but leaves its magnitude,
and hence the Rabi frequency (64), unchanged.

% ============================================================
\section{Superposition of All Propagation Paths}
% ============================================================

For user \(k\), the transmitted signal reaches the vapor cell
through \(L_k\) propagation paths.

The total electric field produced by user \(k\) is therefore the
sum of the fields associated with all its propagation paths:
\[
\mathbf E_{i,k}(t)
=
\sum_{l=1}^{L_k}
\mathbf E_{i,k,l}(t).
\]

Substituting (15),
\[
\mathbf E_{i,k}(t)
=
\sum_{l=1}^{L_k}
\boldsymbol{\epsilon}_{i,k,l}
\alpha_{l,k}s_k
\cos
\left[
\omega t+(i-1)\phi_{l,k}
\right].
\]

The vapor cell receives signals from all \(K\) users. Therefore,
the total electric field at vapor cell \(i\) is
\[
\mathbf E_i(t)
=
\sum_{k=1}^{K}
\mathbf E_{i,k}(t).
\]

Substituting the expression above,
\[
\mathbf E_i(t)
=
\sum_{k=1}^{K}
\sum_{l=1}^{L_k}
\boldsymbol{\epsilon}_{i,k,l}
\alpha_{l,k}s_k
\cos
\left[
\omega t+(i-1)\phi_{l,k}
\right].
\]

Thus,
\begin{equation}
\boxed{
\mathbf E_i(t)
=
\sum_{k=1}^{K}
\sum_{l=1}^{L_k}
\boldsymbol{\epsilon}_{i,k,l}
\alpha_{l,k}s_k
\cos
\left[
\omega t+(i-1)\phi_{l,k}
\right].
}
\tag{16}
\end{equation}

% ============================================================
\section{Electric-Dipole Interaction}
% ============================================================

The electromagnetic field interacts with the electric dipole of
the atom.

For an electric dipole moment \(\boldsymbol{\mu}\) in an electric
field \(\mathbf E\), the interaction energy is represented by the
dipole--field product.

For a charge \(q\) located at position \(\mathbf r\), the electric
dipole moment is
\[
\boldsymbol{\mu}
=
q\mathbf r.
\]

Therefore, the interaction energy can be written in terms of the
position vector and electric field as
\[
V
=
q\mathbf r^{T}\mathbf E.
\]

In quantum mechanics, the position of the electron is represented
by a position operator. Hence,
\[
\hat{\mathbf r}
=
\begin{bmatrix}
\hat x\\
\hat y\\
\hat z
\end{bmatrix}.
\]

The interaction Hamiltonian for the atom in vapor cell \(i\) is
therefore written as
\begin{equation}
\boxed{
\hat V_i
=
q\hat{\mathbf r}^{\,T}\mathbf E_i(t).
}
\tag{17}
\end{equation}

The electric field is a vector,
\[
\mathbf E_i(t)
=
\begin{bmatrix}
E_{i,x}(t)\\
E_{i,y}(t)\\
E_{i,z}(t)
\end{bmatrix},
\]
and hence
\[
\begin{aligned}
\hat{\mathbf r}^{\,T}\mathbf E_i(t)
&=
\begin{bmatrix}
\hat x&\hat y&\hat z
\end{bmatrix}
\begin{bmatrix}
E_{i,x}(t)\\
E_{i,y}(t)\\
E_{i,z}(t)
\end{bmatrix}\\
&=
\hat xE_{i,x}(t)
+
\hat yE_{i,y}(t)
+
\hat zE_{i,z}(t).
\end{aligned}
\]

Therefore,
\begin{equation}
\boxed{
\hat V_i
=
q
\left[
\hat xE_{i,x}(t)
+
\hat yE_{i,y}(t)
+
\hat zE_{i,z}(t)
\right].
}
\tag{18}
\end{equation}

At this stage, the interaction Hamiltonian is expressed in terms
of the position operator. The next step is to represent this
operator in the two-level basis
\[
\{\ketg,\kete\}.
\]

% ============================================================
\section{Two-Level Representation of the Position Operator}
% ============================================================

In quantum mechanics, the position of the electron is represented by
a position operator rather than an ordinary definite coordinate.

The position operator is written as
\[
\hat{\mathbf r}
=
\begin{bmatrix}
\hat x\\
\hat y\\
\hat z
\end{bmatrix}.
\]

For the two-level Rydberg-atom system, the state can be expressed as
\[
\ket{\Phi_i(t)}
=
\zeta_{g,i}(t)\ketg
+
\zeta_{e,i}(t)\kete.
\]

The completeness relation for the two-level basis is
\begin{equation}
\boxed{
\hat I
=
\ketg\bra g
+
\kete\bra e.
}
\tag{19}
\end{equation}

Using the completeness relation on both sides of the position
operator,
\[
\hat{\mathbf r}
=
\hat I\hat{\mathbf r}\hat I.
\]

Therefore,
\[
\hat{\mathbf r}
=
\left(
\ketg\bra g+\kete\bra e
\right)
\hat{\mathbf r}
\left(
\ketg\bra g+\kete\bra e
\right).
\]

Expanding the product gives
\[
\begin{aligned}
\hat{\mathbf r}
={}&
\ketg\bra g
\hat{\mathbf r}
\ketg\bra g
+
\ketg\bra g
\hat{\mathbf r}
\kete\bra e\\
&+
\kete\bra e
\hat{\mathbf r}
\ketg\bra g
+
\kete\bra e
\hat{\mathbf r}
\kete\bra e.
\end{aligned}
\]

The diagonal matrix elements are
\[
\bra g\hat{\mathbf r}\ketg
\]
and
\[
\bra e\hat{\mathbf r}\kete,
\]
which represent the position-operator expectation values within
the ground and excited states.

The off-diagonal matrix elements
\[
\bra e\hat{\mathbf r}\ketg
\]
and
\[
\bra g\hat{\mathbf r}\kete
\]
are the transition matrix elements between the two atomic states.

Atomic eigenstates have definite parity, while $\hat{\mathbf r}$ is an
odd operator. Hence the diagonal elements vanish,
\[
\bra g\hat{\mathbf r}\ketg
=
\bra e\hat{\mathbf r}\kete
=
\mathbf 0,
\]
and only the transition terms remain (the states $\ketg$ and $\kete$
have opposite parity, which is the electric-dipole selection rule).
Hence,
\begin{equation}
\boxed{
\hat{\mathbf r}
=
\bra e\hat{\mathbf r}\ketg
\kete\bra g
+
\bra g\hat{\mathbf r}\kete
\ketg\bra e.
}
\tag{20}
\end{equation}

% ============================================================
\section{Transition Dipole Moment}
% ============================================================

The interaction Hamiltonian obtained previously is
\[
\hat V_i
=
q\hat{\mathbf r}^{\,T}\mathbf E_i(t).
\]

Substituting the two-level representation of the position operator,
\[
\hat V_i
=
q
\left[
\bra e\hat{\mathbf r}\ketg
\kete\bra g
+
\bra g\hat{\mathbf r}\kete
\ketg\bra e
\right]^{T}
\mathbf E_i(t).
\]

The electric field is an ordinary classical vector. Therefore, it
can be included in the corresponding scalar products.

Define the transition dipole moment between the ground and excited
states as
\begin{equation}
\boxed{
\boldsymbol{\mu}_{eg}
=
q\bra e\hat{\mathbf r}\ketg.
}
\tag{21}
\end{equation}

Similarly, define
\begin{equation}
\boxed{
\boldsymbol{\mu}_{ge}
=
q\bra g\hat{\mathbf r}\kete.
}
\tag{22}
\end{equation}

Therefore,
\[
q
\bra e\hat{\mathbf r}\ketg
=
\boldsymbol{\mu}_{eg},
\qquad
q
\bra g\hat{\mathbf r}\kete
=
\boldsymbol{\mu}_{ge}.
\]

Substituting these definitions into the interaction Hamiltonian,
\[
\hat V_i
=
\left(
\boldsymbol{\mu}_{eg}^{T}
\mathbf E_i(t)
\right)
\kete\bra g
+
\left(
\boldsymbol{\mu}_{ge}^{T}
\mathbf E_i(t)
\right)
\ketg\bra e.
\]

The position operator is Hermitian:
\[
\hat{\mathbf r}^{\dagger}
=
\hat{\mathbf r}.
\]

Hence,
\[
\left(
\bra e\hat{\mathbf r}\ketg
\right)^{*}
=
\bra g\hat{\mathbf r}\kete.
\]

Therefore,
\[
\boldsymbol{\mu}_{ge}
=
\boldsymbol{\mu}_{eg}^{*}.
\]

Consequently,
\[
\left(
\boldsymbol{\mu}_{eg}^{T}\mathbf E_i(t)
\right)^{*}
=
\boldsymbol{\mu}_{ge}^{T}\mathbf E_i(t),
\]
for a real electric field.

Define
\begin{equation}
\boxed{
\bar{\omega}_i(t)
=
\boldsymbol{\mu}_{eg}^{T}\mathbf E_i(t).
}
\tag{23}
\end{equation}

Then,
\[
\bar{\omega}_i^{*}(t)
=
\boldsymbol{\mu}_{ge}^{T}\mathbf E_i(t).
\]

Therefore, the interaction Hamiltonian becomes
\begin{equation}
\boxed{
\hat V_i
=
\bar{\omega}_i(t)
\kete\bra g
+
\bar{\omega}_i^{*}(t)
\ketg\bra e.
}
\tag{24}
\end{equation}

% ============================================================
\section{Matrix Representation of the Interaction Hamiltonian}
% ============================================================

Using
\[
\kete
=
\begin{bmatrix}
1\\
0
\end{bmatrix},
\qquad
\ketg
=
\begin{bmatrix}
0\\
1
\end{bmatrix},
\]
we obtain
\[
\kete\bra g
=
\begin{bmatrix}
1\\
0
\end{bmatrix}
\begin{bmatrix}
0&1
\end{bmatrix}
=
\begin{bmatrix}
0&1\\
0&0
\end{bmatrix}.
\]

Similarly,
\[
\ketg\bra e
=
\begin{bmatrix}
0\\
1
\end{bmatrix}
\begin{bmatrix}
1&0
\end{bmatrix}
=
\begin{bmatrix}
0&0\\
1&0
\end{bmatrix}.
\]

Substituting these into (24),
\[
\begin{aligned}
\hat V_i
&=
\bar{\omega}_i(t)
\begin{bmatrix}
0&1\\
0&0
\end{bmatrix}
+
\bar{\omega}_i^{*}(t)
\begin{bmatrix}
0&0\\
1&0
\end{bmatrix}\\
&=
\begin{bmatrix}
0&\bar{\omega}_i(t)\\
\bar{\omega}_i^{*}(t)&0
\end{bmatrix}.
\end{aligned}
\]

Hence,
\begin{equation}
\boxed{
\hat V_i
=
\begin{bmatrix}
0&\bar{\omega}_i(t)\\
\bar{\omega}_i^{*}(t)&0
\end{bmatrix}.
}
\tag{25}
\end{equation}

% ============================================================
\section{Interaction Coefficient in Terms of the Electric Field}
% ============================================================

From the definition
\[
\bar{\omega}_i(t)
=
\boldsymbol{\mu}_{eg}^{T}\mathbf E_i(t),
\]
and the total electric field (16),
\[
\mathbf E_i(t)
=
\sum_{k=1}^{K}
\sum_{l=1}^{L_k}
\boldsymbol{\epsilon}_{i,k,l}
\alpha_{l,k}s_k
\cos
\left[
\omega t+(i-1)\phi_{l,k}
\right],
\]
we obtain
\[
\bar{\omega}_i(t)
=
\boldsymbol{\mu}_{eg}^{T}
\sum_{k=1}^{K}
\sum_{l=1}^{L_k}
\boldsymbol{\epsilon}_{i,k,l}
\alpha_{l,k}s_k
\cos
\left[
\omega t+(i-1)\phi_{l,k}
\right].
\]

Thus,
\begin{equation}
\boxed{
\bar{\omega}_i(t)
=
\sum_{k=1}^{K}
\sum_{l=1}^{L_k}
\boldsymbol{\mu}_{eg}^{T}
\boldsymbol{\epsilon}_{i,k,l}
\alpha_{l,k}s_k
\cos
\left[
\omega t+(i-1)\phi_{l,k}
\right].
}
\tag{26}
\end{equation}

Similarly,
\[
\bar{\omega}_i^{*}(t)
=
\sum_{k=1}^{K}
\sum_{l=1}^{L_k}
\boldsymbol{\mu}_{ge}^{T}
\boldsymbol{\epsilon}_{i,k,l}
\alpha_{l,k}^{*}s_k^{*}
\cos
\left[
\omega t+(i-1)\phi_{l,k}
\right].
\]

% ============================================================
\section{Total Hamiltonian}
% ============================================================

The free Hamiltonian is
\[
\hat H_i
=
\begin{bmatrix}
\hbar\omega_e&0\\
0&\hbar\omega_g
\end{bmatrix},
\]
and the interaction Hamiltonian is
\[
\hat V_i
=
\begin{bmatrix}
0&\bar{\omega}_i(t)\\
\bar{\omega}_i^{*}(t)&0
\end{bmatrix}.
\]

Therefore,
\[
\begin{aligned}
\hat H_i+\hat V_i
&=
\begin{bmatrix}
\hbar\omega_e&0\\
0&\hbar\omega_g
\end{bmatrix}
+
\begin{bmatrix}
0&\bar{\omega}_i(t)\\
\bar{\omega}_i^{*}(t)&0
\end{bmatrix}\\
&=
\begin{bmatrix}
\hbar\omega_e&\bar{\omega}_i(t)\\
\bar{\omega}_i^{*}(t)&\hbar\omega_g
\end{bmatrix}.
\end{aligned}
\]

Hence,
\begin{equation}
\boxed{
\hat H_i+\hat V_i
=
\begin{bmatrix}
\hbar\omega_e&\bar{\omega}_i(t)\\
\bar{\omega}_i^{*}(t)&\hbar\omega_g
\end{bmatrix}.
}
\tag{27}
\end{equation}

% ============================================================
\section{Coupled Schr\"odinger Equations}
% ============================================================

Using
\[
\ket{\Phi_i(t)}
=
\begin{bmatrix}
\zeta_{e,i}(t)\\
\zeta_{g,i}(t)
\end{bmatrix},
\]
the Schr\"odinger equation becomes
\[
i\hbar
\frac{d}{dt}
\begin{bmatrix}
\zeta_{e,i}(t)\\
\zeta_{g,i}(t)
\end{bmatrix}
=
\begin{bmatrix}
\hbar\omega_e&\bar{\omega}_i(t)\\
\bar{\omega}_i^{*}(t)&\hbar\omega_g
\end{bmatrix}
\begin{bmatrix}
\zeta_{e,i}(t)\\
\zeta_{g,i}(t)
\end{bmatrix}.
\]

Multiplying the matrices,
\[
i\hbar
\begin{bmatrix}
\dot{\zeta}_{e,i}(t)\\
\dot{\zeta}_{g,i}(t)
\end{bmatrix}
=
\begin{bmatrix}
\hbar\omega_e\zeta_{e,i}(t)
+
\bar{\omega}_i(t)\zeta_{g,i}(t)
\\[4pt]
\bar{\omega}_i^{*}(t)\zeta_{e,i}(t)
+
\hbar\omega_g\zeta_{g,i}(t)
\end{bmatrix}.
\]

Therefore,
\begin{equation}
\boxed{
i\hbar\dot{\zeta}_{e,i}
=
\hbar\omega_e\zeta_{e,i}
+
\bar{\omega}_i(t)\zeta_{g,i}.
}
\tag{28a}
\end{equation}

and
\begin{equation}
\boxed{
i\hbar\dot{\zeta}_{g,i}
=
\bar{\omega}_i^{*}(t)\zeta_{e,i}
+
\hbar\omega_g\zeta_{g,i}.
}
\tag{28b}
\end{equation}

The terms
\(\hbar\omega_e\zeta_{e,i}\)
and
\(\hbar\omega_g\zeta_{g,i}\)
represent the free atomic evolution, while the terms containing
\(\bar{\omega}_i(t)\) represent the interaction with the
electromagnetic field.

% ============================================================
\section{Removal of the Free-Hamiltonian Evolution}
% ============================================================

We start from the time-dependent Schr\"odinger equation for the
free atomic evolution,
\[
i\hbar
\frac{\partial}{\partial t}
|\Phi_i(t)\rangle
=
\hat H_i|\Phi_i(t)\rangle.
\]

For the free Hamiltonian,
\[
\hat H_i
=
\begin{bmatrix}
\hbar\omega_e & 0\\
0 & \hbar\omega_g
\end{bmatrix}.
\]

Consider first the excited-state component. From the free
Schr\"odinger equation,
\[
i\hbar
\frac{d\zeta_{e,i}(t)}{dt}
=
\hbar\omega_e\zeta_{e,i}(t),
\]
we obtain
\[
\frac{d\zeta_{e,i}(t)}{dt}
=
-j\omega_e\zeta_{e,i}(t).
\]

Now separate the variables:
\[
\frac{1}{\zeta_{e,i}(t)}
\,d\zeta_{e,i}(t)
=
-j\omega_e\,dt.
\]

Integrating both sides,
\[
\int
\frac{1}{\zeta_{e,i}}
\,d\zeta_{e,i}
=
\int
-j\omega_e\,dt,
\]
which gives
\[
\ln\zeta_{e,i}(t)
=
-j\omega_e t+C.
\]

Exponentiating,
\[
\zeta_{e,i}(t)
=
Ce^{-j\omega_e t}.
\]

Thus, the free evolution of the excited-state component contains
the phase factor
\[
\boxed{
e^{-j\omega_e t}
}.
\]

Similarly, for the ground-state component,
\[
i\hbar
\frac{d\zeta_{g,i}(t)}{dt}
=
\hbar\omega_g\zeta_{g,i}(t),
\]
which gives
\[
\frac{d\zeta_{g,i}(t)}{dt}
=
-j\omega_g\zeta_{g,i}(t),
\]
and hence
\[
\boxed{
e^{-j\omega_g t}
}
\]
for the ground state.

We now define the slowly varying amplitudes \(a_i(t)\) and
\(b_i(t)\) by
\begin{equation}
\boxed{
\zeta_{e,i}(t)
=
a_i(t)e^{-j\omega_e t}.
}
\tag{29a}
\end{equation}

Similarly,
\begin{equation}
\boxed{
\zeta_{g,i}(t)
=
b_i(t)e^{-j\omega_g t}.
}
\tag{29b}
\end{equation}

The functions \(a_i(t)\) and \(b_i(t)\) contain the evolution
produced by the electromagnetic interaction after the free atomic
oscillation has been removed.

For the excited state,
\[
\zeta_{e,i}(t)
=
a_i(t)e^{-j\omega_e t}.
\]

Differentiating,
\[
\begin{aligned}
\dot{\zeta}_{e,i}(t)
&=
\frac{d}{dt}
\left[
a_i(t)e^{-j\omega_e t}
\right]\\
&=
\dot a_i(t)e^{-j\omega_e t}
+
a_i(t)
\frac{d}{dt}
\left(
e^{-j\omega_e t}
\right)\\
&=
\dot a_i(t)e^{-j\omega_e t}
-
j\omega_ea_i(t)e^{-j\omega_e t}\\
&=
\left(
\dot a_i-j\omega_ea_i
\right)
e^{-j\omega_e t}.
\end{aligned}
\]

Hence,
\begin{equation}
\boxed{
\dot{\zeta}_{e,i}
=
\left(
\dot a_i-j\omega_ea_i
\right)
e^{-j\omega_e t}.
}
\tag{30}
\end{equation}

Similarly,
\[
\begin{aligned}
\dot{\zeta}_{g,i}
&=
\frac{d}{dt}
\left[
b_i(t)e^{-j\omega_g t}
\right]\\
&=
\dot b_i(t)e^{-j\omega_g t}
-
j\omega_gb_i(t)e^{-j\omega_g t}\\
&=
\left(
\dot b_i-j\omega_gb_i
\right)
e^{-j\omega_g t}.
\end{aligned}
\]

Therefore,
\begin{equation}
\boxed{
\dot{\zeta}_{g,i}
=
\left(
\dot b_i-j\omega_gb_i
\right)
e^{-j\omega_g t}.
}
\tag{31}
\end{equation}

% ============================================================
\section{Transformation of the Excited-State Equation}
% ============================================================

Substitute
\[
\zeta_{e,i}
=
a_i e^{-j\omega_e t},
\qquad
\zeta_{g,i}
=
b_i e^{-j\omega_g t},
\]
and (30) into (28a):
\[
i\hbar
\left(
\dot a_i-j\omega_ea_i
\right)
e^{-j\omega_e t}
=
\hbar\omega_e
a_i e^{-j\omega_e t}
+
\bar{\omega}_i(t)
b_i e^{-j\omega_g t}.
\]

Expanding the left-hand side,
\[
i\hbar\dot a_i e^{-j\omega_e t}
+
\hbar\omega_ea_i e^{-j\omega_e t}
=
\hbar\omega_ea_i e^{-j\omega_e t}
+
\bar{\omega}_i(t)b_i e^{-j\omega_g t}.
\]

The free-Hamiltonian terms cancel:
\[
i\hbar\dot a_i e^{-j\omega_e t}
=
\bar{\omega}_i(t)b_i e^{-j\omega_g t}.
\]

Multiplying both sides by
\(e^{j\omega_e t}\),
we obtain
\[
i\hbar\dot a_i
=
\bar{\omega}_i(t)b_i
e^{j(\omega_e-\omega_g)t}.
\]

Define the atomic transition angular frequency as
\begin{equation}
\boxed{
\omega_{eg}
=
\omega_e-\omega_g.
}
\tag{32}
\end{equation}

Hence,
\begin{equation}
\boxed{
i\hbar\dot a_i
=
\bar{\omega}_i(t)
b_i
e^{j\omega_{eg}t}.
}
\tag{33}
\end{equation}

Similarly, substituting the transformations into the ground-state
equation (28b) gives
\begin{equation}
\boxed{
i\hbar\dot b_i
=
\bar{\omega}_i^{*}(t)
a_i
e^{-j\omega_{eg}t}.
}
\tag{34}
\end{equation}

% ============================================================
\section{Substitution of the Electromagnetic Interaction}
% ============================================================

From (26),
\[
\bar{\omega}_i(t)
=
\sum_{k=1}^{K}
\sum_{l=1}^{L_k}
\boldsymbol{\mu}_{eg}^{T}
\boldsymbol{\epsilon}_{i,k,l}
\alpha_{l,k}s_k
\cos
\left[
\omega t+(i-1)\phi_{l,k}
\right].
\]

Define
\begin{equation}
\boxed{
A_{i,k,l}
=
\boldsymbol{\mu}_{eg}^{T}
\boldsymbol{\epsilon}_{i,k,l}
\alpha_{l,k}s_k.
}
\tag{35}
\end{equation}

Then,
\[
\bar{\omega}_i(t)
=
\sum_{k=1}^{K}
\sum_{l=1}^{L_k}
A_{i,k,l}
\cos
\left[
\omega t+(i-1)\phi_{l,k}
\right].
\]

Substituting this expression into (33),
\[
i\hbar\dot a_i
=
\sum_{k=1}^{K}
\sum_{l=1}^{L_k}
A_{i,k,l}
\cos
\left[
\omega t+(i-1)\phi_{l,k}
\right]
e^{j\omega_{eg}t}
b_i.
\]

Using the identity
\[
\cos x
=
\frac{e^{jx}+e^{-jx}}{2},
\]
we obtain
\[
\cos
\left[
\omega t+(i-1)\phi_{l,k}
\right]
=
\frac{1}{2}
\Big[
e^{j[\omega t+(i-1)\phi_{l,k}]}
+
e^{-j[\omega t+(i-1)\phi_{l,k}]}
\Big].
\]

Therefore,
\[
i\hbar\dot a_i
=
\frac{1}{2}
\sum_{k=1}^{K}
\sum_{l=1}^{L_k}
A_{i,k,l}
\Big[
e^{j[\omega t+(i-1)\phi_{l,k}]}
+
e^{-j[\omega t+(i-1)\phi_{l,k}]}
\Big]
e^{j\omega_{eg}t}
b_i.
\]

Combining the exponential terms,
\begin{equation}
\begin{aligned}
i\hbar\dot a_i
=
\frac{1}{2}
\sum_{k=1}^{K}
\sum_{l=1}^{L_k}
A_{i,k,l}
\Big[
&
e^{j(\omega+\omega_{eg})t}
e^{j(i-1)\phi_{l,k}}
\\
&+
e^{j(\omega_{eg}-\omega)t}
e^{-j(i-1)\phi_{l,k}}
\Big]
b_i.
\end{aligned}
\tag{36}
\end{equation}

% ============================================================
\section{Resonance Condition}
% ============================================================

The electromagnetic field is considered resonant with the atomic
transition when
\begin{equation}
\boxed{
\omega=\omega_{eg}.
}
\tag{37}
\end{equation}

Therefore,
\[
\omega+\omega_{eg}=2\omega,
\qquad
\omega_{eg}-\omega=0.
\]

Substituting these relations into (36),
\begin{equation}
i\hbar\dot a_i
=
\frac{1}{2}
\sum_{k=1}^{K}
\sum_{l=1}^{L_k}
A_{i,k,l}
\Big[
e^{j2\omega t}
e^{j(i-1)\phi_{l,k}}
+
e^{-j(i-1)\phi_{l,k}}
\Big]
b_i.
\tag{38}
\end{equation}

% ============================================================
\section{Rotating-Wave Approximation}
% ============================================================

Under the rotating-wave approximation, terms oscillating rapidly
compared with the atomic transition dynamics are neglected.

The first term in (38),
\[
e^{j2\omega t}
e^{j(i-1)\phi_{l,k}},
\]
contains the rapidly oscillating factor
\(e^{j2\omega t}\).

The second term,
\[
e^{-j(i-1)\phi_{l,k}},
\]
is time-independent.

Therefore, the rapidly oscillating term is neglected, giving
\begin{equation}
\boxed{
i\hbar\dot a_i
=
\frac{1}{2}
\sum_{k=1}^{K}
\sum_{l=1}^{L_k}
A_{i,k,l}
e^{-j(i-1)\phi_{l,k}}
b_i.
}
\tag{39}
\end{equation}

Substituting the definition of \(A_{i,k,l}\),
\[
i\hbar\dot a_i
=
\frac{1}{2}
\sum_{k=1}^{K}
\sum_{l=1}^{L_k}
\boldsymbol{\mu}_{eg}^{T}
\boldsymbol{\epsilon}_{i,k,l}
\alpha_{l,k}s_k
e^{-j(i-1)\phi_{l,k}}
b_i.
\]

Define
\begin{equation}
\boxed{
C_i
=
\sum_{k=1}^{K}
\sum_{l=1}^{L_k}
\boldsymbol{\mu}_{eg}^{T}
\boldsymbol{\epsilon}_{i,k,l}
\alpha_{l,k}s_k
e^{-j(i-1)\phi_{l,k}}.
}
\tag{40}
\end{equation}

Then,
\begin{equation}
\boxed{
i\hbar\dot a_i
=
\frac{C_i}{2}b_i.
}
\tag{41}
\end{equation}

Similarly, substituting $\bar\omega_i^*(t)$ into (34) and neglecting
the $e^{-j2\omega t}$ term,
\begin{equation}
\boxed{
i\hbar\dot b_i
=
\frac{C_i^{*}}{2}a_i.
}
\tag{42}
\end{equation}

% ============================================================
\section{Second-Order Differential Equation}
% ============================================================

From (41),
\[
i\hbar\dot a_i
=
\frac{C_i}{2}b_i.
\]

Since \(C_i\) is constant with respect to time in the resonant
RWA model,
\[
i\hbar\ddot a_i
=
\frac{C_i}{2}\dot b_i.
\]

From (42),
\[
\dot b_i
=
\frac{C_i^{*}}{2i\hbar}a_i.
\]

Using
\(1/i=-j\),
we obtain
\[
\dot b_i
=
-\frac{jC_i^{*}}{2\hbar}a_i.
\]

Substituting this into the differentiated form of (41),
\[
\begin{aligned}
i\hbar\ddot a_i
&=
\frac{C_i}{2}
\left(
-\frac{jC_i^{*}}{2\hbar}a_i
\right)
=
-\frac{jC_iC_i^{*}}{4\hbar}a_i
=
-\frac{j|C_i|^2}{4\hbar}a_i.
\end{aligned}
\]

Dividing by \(i\hbar\),
\[
\ddot a_i
=
-\frac{|C_i|^2}{4\hbar^2}a_i.
\]

Therefore,
\begin{equation}
\boxed{
\ddot a_i
+
\frac{|C_i|^2}{4\hbar^2}a_i
=
0.
}
\tag{43}
\end{equation}

% ============================================================
\section{Solution of the Excited-State Dynamics}
% ============================================================

The second-order differential equation obtained in (43) is
\begin{equation}
\ddot{a}_i(t)
+
\frac{|C_i|^2}{4\hbar^2}a_i(t)
=0.
\tag{44}
\end{equation}

Its general solution is
\begin{equation}
a_i(t)
=
A\cos\left(\frac{|C_i|t}{2\hbar}\right)
+
B\sin\left(\frac{|C_i|t}{2\hbar}\right).
\tag{45}
\end{equation}

The atom is initially prepared in the ground state. Therefore,
\begin{equation}
\zeta_{e,i}(0)=0,
\qquad
\zeta_{g,i}(0)=1.
\tag{46}
\end{equation}

Since
\begin{equation}
\zeta_{e,i}(t)=a_i(t)e^{-j\omega_e t},
\tag{47}
\end{equation}
and $e^{-j\omega_e 0}=1$, the initial excited-state condition gives
\begin{equation}
a_i(0)=0.
\tag{48}
\end{equation}

Evaluating (45) at $t=0$ gives
\begin{equation}
a_i(0)=A.
\tag{49}
\end{equation}

Hence,
\begin{equation}
A=0.
\tag{50}
\end{equation}

Therefore, the transformed excited-state amplitude becomes
\begin{equation}
a_i(t)
=
B\sin\left(\frac{|C_i|t}{2\hbar}\right).
\tag{51}
\end{equation}

The coefficient $B$ is determined from the first-order equation
obtained from the rotating-wave approximation,
\begin{equation}
i\hbar\dot a_i(t)=\frac{C_i}{2}b_i(t).
\tag{52}
\end{equation}

At $t=0$, $b_i(0)=\zeta_{g,i}(0)=1$, and therefore
\begin{equation}
i\hbar\dot a_i(0)=\frac{C_i}{2}.
\tag{53}
\end{equation}

Differentiating (51) gives
\begin{equation}
\dot a_i(t)
=
B\frac{|C_i|}{2\hbar}
\cos\left(\frac{|C_i|t}{2\hbar}\right).
\tag{54}
\end{equation}

Thus, at $t=0$,
\begin{equation}
\dot a_i(0)=B\frac{|C_i|}{2\hbar}.
\tag{55}
\end{equation}

Substituting (55) into (53) gives
\begin{equation}
iB|C_i|=C_i.
\tag{56}
\end{equation}

For $C_i\neq0$,
\begin{equation}
B=-j\frac{C_i}{|C_i|}.
\tag{57}
\end{equation}

Hence,
\begin{equation}
a_i(t)
=
-j\frac{C_i}{|C_i|}
\sin\left(\frac{|C_i|t}{2\hbar}\right).
\tag{58}
\end{equation}

Using (47), the original excited-state amplitude is
\begin{equation}
\zeta_{e,i}(t)
=
-j\frac{C_i}{|C_i|}
\sin\left(\frac{|C_i|t}{2\hbar}\right)e^{-j\omega_e t}.
\tag{59}
\end{equation}

The excited-state probability is defined as the squared magnitude of
the excited-state probability amplitude,
\begin{equation}
P_e(t)=|\zeta_{e,i}(t)|^2.
\tag{60}
\end{equation}

Using the multiplicative property of the complex magnitude,
\begin{equation}
|abc|^2=|a|^2|b|^2|c|^2,
\tag{61}
\end{equation}
and the identities
\begin{equation}
|-j|^2=1,
\qquad
\left|\frac{C_i}{|C_i|}\right|^2=1,
\qquad
|e^{-j\omega_e t}|^2=1,
\tag{62}
\end{equation}
we obtain
\begin{equation}
P_e(t)
=
\sin^2\left(\frac{|C_i|t}{2\hbar}\right).
\tag{63}
\end{equation}

Define the Rabi frequency by
\begin{equation}
\boxed{
\Omega_i=\frac{|C_i|}{\hbar}
}.
\tag{64}
\end{equation}

Therefore,
\begin{equation}
\frac{|C_i|t}{2\hbar}=\frac{\Omega_i t}{2}.
\tag{65}
\end{equation}

Substituting (65) into (63) gives the final excited-state probability,
\begin{equation}
\boxed{
P_e(t)=\sin^2\left(\frac{\Omega_i t}{2}\right)
}.
\tag{66}
\end{equation}
Equivalently,
\begin{equation}
\boxed{
|\zeta_{e,i}(t)|^2=\sin^2\left(\frac{\Omega_i t}{2}\right)
}.
\tag{66a}
\end{equation}

% ============================================================
\section{Expression for the Rabi Frequency}
% ============================================================

From the definition of the effective interaction coefficient,
\begin{equation}
C_i
=
\sum_{k=1}^{K}
\sum_{l=1}^{L_k}
\boldsymbol{\mu}_{eg}^{T}
\boldsymbol{\epsilon}_{i,k,l}
\alpha_{l,k}s_k
 e^{-j(i-1)\phi_{l,k}}.
\tag{67}
\end{equation}

Using the definition of the Rabi frequency in (64),
\begin{equation}
\Omega_i
=
\frac{1}{\hbar}
\left|
\sum_{k=1}^{K}
\sum_{l=1}^{L_k}
\boldsymbol{\mu}_{eg}^{T}
\boldsymbol{\epsilon}_{i,k,l}
\alpha_{l,k}s_k
 e^{-j(i-1)\phi_{l,k}}
\right|.
\tag{68}
\end{equation}

\textbf{Remark.} Xu \emph{et al.} write Eq.~(6) with
$\boldsymbol\mu_{eg}^{H}$ and $e^{+j(i-1)\phi_{l,k}}$, but the channel
Eq.~(7) with $\boldsymbol\mu_{eg}^{T}$ and $e^{-j(i-1)\phi_{l,k}}$.
Since $\boldsymbol\mu_{eg}=[0,\,1785.9\,qa_0,\,0]^T$ and
$\boldsymbol\epsilon_{i,k,l}$ are real, the two sums are complex
conjugates of each other up to the conjugation of $\alpha_{l,k}s_k$,
and their magnitudes are equal. Here the form consistent with
Eq.~(7) is used.

Since $s_k$ does not depend on the path index $l$, it can be factored
outside the inner summation:
\begin{equation}
\Omega_i
=
\frac{1}{\hbar}
\left|
\sum_{k=1}^{K}
 s_k
\sum_{l=1}^{L_k}
\boldsymbol{\mu}_{eg}^{T}
\boldsymbol{\epsilon}_{i,k,l}
\alpha_{l,k}
 e^{-j(i-1)\phi_{l,k}}
\right|.
\tag{69}
\end{equation}

Equivalently,
\begin{equation}
\Omega_i
=
\left|
\sum_{k=1}^{K}
 s_k
\left(
\sum_{l=1}^{L_k}
\frac{1}{\hbar}
\boldsymbol{\mu}_{eg}^{T}
\boldsymbol{\epsilon}_{i,k,l}
\alpha_{l,k}
 e^{-j(i-1)\phi_{l,k}}
\right)
\right|.
\tag{70}
\end{equation}

Define the channel coefficient
\begin{equation}
\boxed{
g_{i,k}
=
\sum_{l=1}^{L_k}
\frac{1}{\hbar}
\boldsymbol{\mu}_{eg}^{T}
\boldsymbol{\epsilon}_{i,k,l}
\alpha_{l,k}
e^{-j(i-1)\phi_{l,k}}
}.
\tag{71}
\end{equation}

Therefore,
\begin{equation}
\boxed{
\Omega_i
=
\left|
\sum_{k=1}^{K}
g_{i,k}s_k
\right|
}.
\tag{72}
\end{equation}

% ============================================================
\section{Received Signal Model}
% ============================================================

For the $p$-th time slot, the measured signal at vapor cell $i$ is
modeled as
\begin{equation}
\boxed{
y_{i,p}
=
\left|
\sum_{k=1}^{K}
g_{i,k}s_{k,p}
+
b_{i,p}
+
n_{i,p}
\right|
}.
\tag{73}
\end{equation}

For a single communication user $k$, the contribution received at
vapor cell $i$ during time slot $p$ is
\begin{equation}
g_{i,k}s_{k,p}.
\tag{74}
\end{equation}

Thus, the channel coefficient represents the complex coupling
between the transmitted symbol and the received contribution:
\begin{equation}
\text{received contribution}
=
\text{channel coefficient}
\times
\text{transmitted symbol}.
\tag{75}
\end{equation}

For $K$ communication users, the contributions arriving at the same
vapor cell are superposed linearly. Hence,
\begin{equation}
a_{i,p}
=
\sum_{k=1}^{K}
g_{i,k}s_{k,p},
\tag{76}
\end{equation}
where $a_{i,p}$ denotes the total desired wireless signal at vapor
cell $i$ and time slot $p$.

The channel coefficient in (71) is obtained from the atomic coupling
term,
\begin{equation}
\frac{1}{\hbar}
\boldsymbol{\mu}_{eg}^{T}
\boldsymbol{\epsilon}_{i,k,l}
\alpha_{l,k}
e^{-j(i-1)\phi_{l,k}},
\tag{77}
\end{equation}
which contains the atomic transition dipole coupling, polarization,
path amplitude, and spatial propagation phase.

% ============================================================
\section{Reference Signal}
% ============================================================

In addition to the communication signals, a reference source is
introduced.

Let the reference source transmit the symbol
\begin{equation}
s_{b,p},
\tag{78}
\end{equation}
where the subscript $b$ denotes the reference beam.

The reference source produces an electromagnetic field at vapor cell
$i$. Its contribution is therefore written in the same form as a
single propagation path:
\begin{equation}
b_{i,p}
=
g_{b,i}s_{b,p},
\tag{79}
\end{equation}
where $g_{b,i}$ represents the channel coupling from the reference
source to vapor cell $i$.

For the communication users, multiple propagation paths are
considered:
\begin{equation}
g_{i,k}
=
\sum_{l=1}^{L_k}
\left(
\text{contribution of path }l
\right).
\tag{80}
\end{equation}

For the reference source, the model uses one effective propagation
path. Therefore, the reference channel contains one effective set of
parameters,
\begin{equation}
\boldsymbol{\epsilon}_{b,i},
\qquad
\alpha_b,
\qquad
\phi_b.
\tag{81}
\end{equation}

Starting from the same single-path channel structure,
\begin{equation}
g_{i,k}^{(\mathrm{single})}
=
\frac{1}{\hbar}
\boldsymbol{\mu}_{eg}^{T}
\boldsymbol{\epsilon}_{i,k}
\alpha_k
e^{-j(i-1)\phi_k},
\tag{82}
\end{equation}
the reference channel is obtained by replacing the generic path
parameters with the reference parameters:
\begin{equation}
\boldsymbol{\epsilon}_{i,k}
\rightarrow
\boldsymbol{\epsilon}_{b,i},
\qquad
\alpha_k
\rightarrow
\alpha_b,
\qquad
\phi_k
\rightarrow
\phi_b.
\tag{83}
\end{equation}

Hence,
\begin{equation}
\boxed{
g_{b,i}
=
\frac{1}{\hbar}
\boldsymbol{\mu}_{eg}^{T}
\boldsymbol{\epsilon}_{b,i}
\alpha_b
e^{-j(i-1)\phi_b}
}.
\tag{84}
\end{equation}

Substituting (84) into (79),
\begin{equation}
\boxed{
b_{i,p}
=
s_{b,p}
\left[
\frac{1}{\hbar}
\boldsymbol{\mu}_{eg}^{T}
\boldsymbol{\epsilon}_{b,i}
\alpha_b
e^{-j(i-1)\phi_b}
\right]
}.
\tag{85}
\end{equation}

\textbf{Assumption.} The linear phase progression $(i-1)\phi_b$ treats
the reference (LO) field as a plane wave across the array, as in
Eq.~(9) of Xu \emph{et al.}

% ============================================================
\section{Complete Received Signal}
% ============================================================

The receiver noise at vapor cell $i$ and time slot $p$ is denoted by
\begin{equation}
n_{i,p}
\sim
\mathcal{CN}(0,\sigma^2),
\tag{86}
\end{equation}
modelled as additive in the Rabi-frequency domain and i.i.d.\ over
$(i,p)$.

The total complex quantity experienced by the atomic receiver is
therefore
\begin{equation}
a_{i,p}
+
b_{i,p}
+
n_{i,p},
\tag{87}
\end{equation}
or, using (76),
\begin{equation}
\sum_{k=1}^{K}
g_{i,k}s_{k,p}
+
b_{i,p}
+
n_{i,p}.
\tag{88}
\end{equation}

The Rydberg atomic receiver provides a magnitude-based measurement
of this total coupling. Consequently,
\begin{equation}
\boxed{
y_{i,p}
=
\left|
a_{i,p}
+
b_{i,p}
+
n_{i,p}
\right|
}.
\tag{89}
\end{equation}

Using (76),
\begin{equation}
\boxed{
y_{i,p}
=
\left|
\sum_{k=1}^{K}
g_{i,k}s_{k,p}
+
b_{i,p}
+
n_{i,p}
\right|
}.
\tag{90}
\end{equation}

% ============================================================
\section{Strong-Reference Approximation}
% ============================================================

Define
\begin{equation}
x_{i,p}
=
a_{i,p}+n_{i,p}.
\tag{91}
\end{equation}

Then (89) becomes
\begin{equation}
y_{i,p}
=
|b_{i,p}+x_{i,p}|.
\tag{92}
\end{equation}

Squaring both sides,
\begin{equation}
y_{i,p}^{\,2}
=
|b_{i,p}+x_{i,p}|^2.
\tag{93}
\end{equation}

For any complex quantity $z$,
\begin{equation}
|z|^2=zz^*.
\tag{94}
\end{equation}

Therefore,
\begin{equation}
y_{i,p}^{\,2}
=
(b_{i,p}+x_{i,p})
(b_{i,p}+x_{i,p})^*.
\tag{95}
\end{equation}

Using
\begin{equation}
(b+x)^*=b^*+x^*,
\tag{96}
\end{equation}
gives
\begin{align}
y_{i,p}^{\,2}
&=
(b_{i,p}+x_{i,p})
(b_{i,p}^*+x_{i,p}^*)
\nonumber\\
&=
b_{i,p}b_{i,p}^*
+
b_{i,p}x_{i,p}^*
+
x_{i,p}b_{i,p}^*
+
x_{i,p}x_{i,p}^*.
\tag{97}
\end{align}

Since
\begin{equation}
|b_{i,p}|^2=b_{i,p}b_{i,p}^*,
\qquad
|x_{i,p}|^2=x_{i,p}x_{i,p}^*,
\tag{98}
\end{equation}
and
\begin{equation}
z+z^*=2\operatorname{Re}(z),
\tag{99}
\end{equation}
the cross terms satisfy
\begin{equation}
b_{i,p}x_{i,p}^*
+
b_{i,p}^*x_{i,p}
=
2\operatorname{Re}
\left(
b_{i,p}^*x_{i,p}
\right).
\tag{100}
\end{equation}

Hence,
\begin{equation}
\boxed{
y_{i,p}^{\,2}
=
|b_{i,p}|^2
+
|x_{i,p}|^2
+
2\operatorname{Re}
\left(
b_{i,p}^*x_{i,p}
\right)
}.
\tag{101}
\end{equation}

Taking the square root,
\begin{equation}
y_{i,p}
=
\sqrt{
|b_{i,p}|^2
+
|x_{i,p}|^2
+
2\operatorname{Re}
\left(
b_{i,p}^*x_{i,p}
\right)
}.
\tag{102}
\end{equation}

Factor $|b_{i,p}|^2$ from the square root:
\begin{equation}
y_{i,p}
=
|b_{i,p}|
\sqrt{
1+
\frac{|x_{i,p}|^2}{|b_{i,p}|^2}
+
\frac{
2\operatorname{Re}(b_{i,p}^*x_{i,p})
}{
|b_{i,p}|^2
}
}.
\tag{103}
\end{equation}

Using
\begin{equation}
|b_{i,p}|^2=b_{i,p}b_{i,p}^*,
\tag{104}
\end{equation}
the last term becomes
\begin{equation}
\frac{
b_{i,p}^*x_{i,p}
}{
b_{i,p}b_{i,p}^*
}
=
\frac{x_{i,p}}{b_{i,p}}.
\tag{105}
\end{equation}

Thus,
\begin{equation}
\boxed{
y_{i,p}
=
|b_{i,p}|
\sqrt{
1+
\frac{|x_{i,p}|^2}{|b_{i,p}|^2}
+
2\operatorname{Re}
\left(
\frac{x_{i,p}}{b_{i,p}}
\right)
}
}.
\tag{106}
\end{equation}

% ============================================================
\section{Taylor Expansion Under Strong Reference}
% ============================================================

The strong-reference assumption is
\begin{equation}
|b_{i,p}|
\gg
|x_{i,p}|
=
|a_{i,p}+n_{i,p}|.
\tag{107}
\end{equation}

Consequently,
\begin{equation}
\left|
\frac{x_{i,p}}{b_{i,p}}
\right|
\ll 1.
\tag{108}
\end{equation}

Define
\begin{equation}
r_{i,p}
=
\frac{x_{i,p}}{b_{i,p}}.
\tag{109}
\end{equation}

Then
\begin{equation}
|r_{i,p}|\ll1,
\tag{110}
\end{equation}
and (106) becomes
\begin{equation}
y_{i,p}
=
|b_{i,p}|
\sqrt{
1+
|r_{i,p}|^2
+
2\operatorname{Re}(r_{i,p})
}.
\tag{111}
\end{equation}

Let
\begin{equation}
u
=
|r_{i,p}|^2
+
2\operatorname{Re}(r_{i,p}).
\tag{112}
\end{equation}

For $|u|\ll1$, the first-order Taylor expansion is
\begin{equation}
\sqrt{1+u}
\approx
1+\frac{u}{2}.
\tag{113}
\end{equation}

Then,
\begin{equation}
y_{i,p}
\approx
|b_{i,p}|
\left[
1+
\frac{1}{2}
\left(
|r_{i,p}|^2
+
2\operatorname{Re}(r_{i,p})
\right)
\right].
\tag{114}
\end{equation}

Therefore,
\begin{equation}
y_{i,p}
\approx
|b_{i,p}|
\left[
1+
\frac{1}{2}|r_{i,p}|^2
+
\operatorname{Re}(r_{i,p})
\right].
\tag{115}
\end{equation}

Substituting (109),
\begin{equation}
y_{i,p}
\approx
|b_{i,p}|
\left[
1+
\frac{1}{2}
\frac{|x_{i,p}|^2}{|b_{i,p}|^2}
+
\operatorname{Re}
\left(
\frac{x_{i,p}}{b_{i,p}}
\right)
\right],
\tag{116}
\end{equation}
which is Eq.~(10) of Xu \emph{et al.}

Expanding the factor $|b_{i,p}|$,
\begin{equation}
y_{i,p}
\approx
|b_{i,p}|
+
\frac{|x_{i,p}|^2}{2|b_{i,p}|}
+
|b_{i,p}|
\operatorname{Re}
\left(
\frac{x_{i,p}}{b_{i,p}}
\right).
\tag{117}
\end{equation}

The second term is second order in the small quantity
$x_{i,p}/b_{i,p}$, whereas the third term is first order. Under the
strong-reference approximation, the second-order term is neglected.

Hence,
\begin{equation}
\boxed{
y_{i,p}
\approx
|b_{i,p}|
+
|b_{i,p}|
\operatorname{Re}
\left(
\frac{x_{i,p}}{b_{i,p}}
\right)
}.
\tag{118}
\end{equation}

Since
\begin{equation}
x_{i,p}
=
a_{i,p}+n_{i,p},
\tag{119}
\end{equation}
we obtain
\begin{equation}
y_{i,p}
\approx
|b_{i,p}|
+
|b_{i,p}|
\operatorname{Re}
\left(
\frac{a_{i,p}+n_{i,p}}{b_{i,p}}
\right).
\tag{120}
\end{equation}

% ============================================================
\section{Reference Magnitude and Phase Representation}
% ============================================================

Represent the reference signal in magnitude-phase form as
\begin{equation}
b_{i,p}
=
|b_{i,p}|e^{jz_{i,p}},
\tag{121}
\end{equation}
where
\begin{equation}
z_{i,p}
=
\angle b_{i,p}
=
\arg(b_{i,p}).
\tag{122}
\end{equation}

Therefore,
\begin{equation}
\frac{|b_{i,p}|}{b_{i,p}}
=
e^{-jz_{i,p}}.
\tag{123}
\end{equation}

Since $|b_{i,p}|$ is real, it can be moved inside the real part of
(120). Using (123),
\begin{equation}
y_{i,p}
\approx
|b_{i,p}|
+
\operatorname{Re}
\left(
e^{-jz_{i,p}}
(a_{i,p}+n_{i,p})
\right).
\tag{124}
\end{equation}

Using the linearity of the real-part operator,
\begin{equation}
\operatorname{Re}(A+B)
=
\operatorname{Re}(A)
+
\operatorname{Re}(B),
\tag{125}
\end{equation}
gives
\begin{equation}
y_{i,p}
\approx
|b_{i,p}|
+
\operatorname{Re}
\left(
e^{-jz_{i,p}}a_{i,p}
\right)
+
\operatorname{Re}
\left(
e^{-jz_{i,p}}n_{i,p}
\right).
\tag{126}
\end{equation}

Define the effective real-valued noise term
\begin{equation}
\bar n_{i,p}
=
\operatorname{Re}
\left(
e^{-jz_{i,p}}n_{i,p}
\right).
\tag{127}
\end{equation}

Since $n_{i,p}$ is circularly symmetric, $e^{-jz_{i,p}}n_{i,p}\sim
\mathcal{CN}(0,\sigma^2)$, and its real part carries half of the
variance:
\begin{equation}
\boxed{
\bar n_{i,p}
\sim
\mathcal{N}\!\left(0,\frac{\sigma^2}{2}\right)
}.
\tag{127a}
\end{equation}

Thus,
\begin{equation}
\boxed{
y_{i,p}
=
\operatorname{Re}
\left(
e^{-jz_{i,p}}a_{i,p}
\right)
+
|b_{i,p}|
+
\bar n_{i,p}
}.
\tag{128}
\end{equation}

Finally, using
\begin{equation}
a_{i,p}
=
\sum_{k=1}^{K}
g_{i,k}s_{k,p},
\tag{129}
\end{equation}
the linearized measurement model becomes
\begin{equation}
\boxed{
y_{i,p}
=
\operatorname{Re}
\left(
e^{-jz_{i,p}}
\sum_{k=1}^{K}
g_{i,k}s_{k,p}
\right)
+
|b_{i,p}|
+
\bar n_{i,p}
}.
\tag{130}
\end{equation}

% ============================================================
\section{Matrix Representation of the Linearized Measurement Model}
% ============================================================

From the scalar linearized measurement model obtained previously,
\begin{equation}
y_{i,p}
\approx
\operatorname{Re}
\left(
e^{-jz_{i,p}}
\sum_{k=1}^{K}
g_{i,k}s_{k,p}
\right)
+
|b_{i,p}|
+
\bar n_{i,p}.
\tag{131}
\end{equation}

The indices are defined as
\[
i=1,2,\ldots,I
\]
for the $I$ vapor cells or receiving antennas,
\[
k=1,2,\ldots,K
\]
for the $K$ users, and
\[
p=1,2,\ldots,P
\]
for the $P$ transmitted symbols or time slots.

Thus, $y_{i,p}$ represents one scalar magnitude measurement
corresponding to vapor cell $i$ and time slot $p$.

Collecting all measurements gives the matrix
\begin{equation}
\boxed{
Y=[y_{i,p}]
\in
\mathbb{R}^{I\times P}
}.
\tag{132}
\end{equation}

Explicitly,
\begin{equation}
Y=
\begin{bmatrix}
y_{1,1} & y_{1,2} & \cdots & y_{1,P}\\
y_{2,1} & y_{2,2} & \cdots & y_{2,P}\\
\vdots & \vdots & \ddots & \vdots\\
y_{I,1} & y_{I,2} & \cdots & y_{I,P}
\end{bmatrix}.
\tag{133}
\end{equation}

Since $y_{i,p}$ is a magnitude measurement,
\begin{equation}
y_{i,p}\in\mathbb{R}.
\tag{134}
\end{equation}

\subsection{Construction of the Phase Matrix}

The scalar model contains the phase factor
\begin{equation}
e^{-jz_{i,p}}.
\tag{135}
\end{equation}

Define the phase matrix
\begin{equation}
\boxed{
Z=
[e^{-jz_{i,p}}]
\in
\mathbb{C}^{I\times P}
}.
\tag{136}
\end{equation}

The matrix is therefore
\begin{equation}
Z=
\begin{bmatrix}
e^{-jz_{1,1}} & e^{-jz_{1,2}} & \cdots & e^{-jz_{1,P}}\\
e^{-jz_{2,1}} & e^{-jz_{2,2}} & \cdots & e^{-jz_{2,P}}\\
\vdots & \vdots & \ddots & \vdots\\
e^{-jz_{I,1}} & e^{-jz_{I,2}} & \cdots & e^{-jz_{I,P}}
\end{bmatrix}.
\tag{137}
\end{equation}

The phase variable is determined from the reference signal according
to
\begin{equation}
z_{i,p}
=
\angle b_{i,p}.
\tag{138}
\end{equation}

\subsection{Construction of the Channel Matrix}

The channel coefficient between vapor cell $i$ and user $k$ is
$g_{i,k}$.

There are $I$ receiving cells and $K$ users. Hence, the channel
matrix is
\begin{equation}
\boxed{
G=[g_{i,k}]
\in
\mathbb{C}^{I\times K}
}.
\tag{139}
\end{equation}

Explicitly,
\begin{equation}
G=
\begin{bmatrix}
g_{1,1} & g_{1,2} & \cdots & g_{1,K}\\
g_{2,1} & g_{2,2} & \cdots & g_{2,K}\\
\vdots & \vdots & \ddots & \vdots\\
g_{I,1} & g_{I,2} & \cdots & g_{I,K}
\end{bmatrix}.
\tag{140}
\end{equation}

Similarly, the transmitted-symbol matrix is defined as
\begin{equation}
\boxed{
S=[s_{k,p}]
\in
\mathbb{C}^{K\times P}
}.
\tag{141}
\end{equation}

Explicitly,
\begin{equation}
S=
\begin{bmatrix}
s_{1,1} & s_{1,2} & \cdots & s_{1,P}\\
s_{2,1} & s_{2,2} & \cdots & s_{2,P}\\
\vdots & \vdots & \ddots & \vdots\\
s_{K,1} & s_{K,2} & \cdots & s_{K,P}
\end{bmatrix}.
\tag{142}
\end{equation}

The matrix product $GS$ has dimensions
\begin{equation}
G S:
\qquad
(I\times K)(K\times P)
=
I\times P.
\tag{143}
\end{equation}

Its $(i,p)$-th element is
\begin{equation}
(GS)_{i,p}
=
\sum_{k=1}^{K}
g_{i,k}s_{k,p}.
\tag{144}
\end{equation}

Therefore,
\begin{equation}
(GS)_{i,p}=a_{i,p}.
\tag{145}
\end{equation}

\subsection{Hadamard Product with the Phase Matrix}

The phase factor $e^{-jz_{i,p}}$ has to multiply the corresponding
$(i,p)$-th element of $GS$.

For two matrices $A$ and $B$ having the same dimensions, the Hadamard
product is defined element-wise as
\begin{equation}
(A\circ B)_{i,p}
=
A_{i,p}B_{i,p}.
\tag{146}
\end{equation}

Therefore,
\begin{equation}
[Z\circ(GS)]_{i,p}
=
Z_{i,p}(GS)_{i,p}.
\tag{147}
\end{equation}

Using
\begin{equation}
Z_{i,p}
=
e^{-jz_{i,p}},
\tag{148}
\end{equation}
\begin{equation}
[Z\circ(GS)]_{i,p}
=
e^{-jz_{i,p}}
\sum_{k=1}^{K}
g_{i,k}s_{k,p}.
\tag{149}
\end{equation}

Applying the real-part operator gives
\begin{equation}
\left(
\operatorname{Re}[Z\circ(GS)]
\right)_{i,p}
=
\operatorname{Re}
\left(
e^{-jz_{i,p}}
\sum_{k=1}^{K}
g_{i,k}s_{k,p}
\right).
\tag{150}
\end{equation}

\subsection{Construction of the Reference Magnitude Matrix}

The scalar measurement model contains the reference magnitude
\begin{equation}
|b_{i,p}|.
\tag{151}
\end{equation}

Define the reference matrix as
\begin{equation}
\boxed{
B=[b_{i,p}]
\in
\mathbb{C}^{I\times P}
}.
\tag{152}
\end{equation}

The corresponding element-wise absolute-value matrix is
\begin{equation}
\boxed{
|B|=[|b_{i,p}|]
\in
\mathbb{R}^{I\times P}
}.
\tag{153}
\end{equation}

Here $|B|$ denotes the element-wise absolute value and not a matrix
norm. Thus,
\begin{equation}
|B|_{i,p}
=
|b_{i,p}|.
\tag{154}
\end{equation}

\subsection{Construction of the Noise Matrix}

The effective real-valued noise term is $\bar n_{i,p}$. Define
\begin{equation}
\boxed{
N=[\bar n_{i,p}]
\in
\mathbb{R}^{I\times P}
}.
\tag{155}
\end{equation}

Explicitly,
\begin{equation}
N=
\begin{bmatrix}
\bar n_{1,1} & \bar n_{1,2} & \cdots & \bar n_{1,P}\\
\bar n_{2,1} & \bar n_{2,2} & \cdots & \bar n_{2,P}\\
\vdots & \vdots & \ddots & \vdots\\
\bar n_{I,1} & \bar n_{I,2} & \cdots & \bar n_{I,P}
\end{bmatrix}.
\tag{156}
\end{equation}

% ============================================================
\section{Complete Matrix Measurement Model}
% ============================================================

The scalar linearized model is
\begin{equation}
y_{i,p}
=
\operatorname{Re}
\left(
e^{-jz_{i,p}}
\sum_{k=1}^{K}
g_{i,k}s_{k,p}
\right)
+
|b_{i,p}|
+
\bar n_{i,p}.
\tag{157}
\end{equation}

Using the definitions above, the element-wise representation becomes
\begin{equation}
y_{i,p}
=
\operatorname{Re}
\left(
[Z\circ(GS)]_{i,p}
\right)
+
|B|_{i,p}
+
N_{i,p}.
\tag{158}
\end{equation}

Collecting all elements into matrix form gives
\begin{equation}
\boxed{
Y
=
\operatorname{Re}[Z\circ(GS)]
+
|B|
+
N
}.
\tag{159}
\end{equation}

The Hadamard product is commutative,
\begin{equation}
Z\circ(GS)
=
(GS)\circ Z.
\tag{160}
\end{equation}

Therefore, either ordering produces the same element-wise product,
and (159) is Eq.~(12) of Xu \emph{et al.}

The use of the Hadamard product is required because each phase
coefficient $e^{-jz_{i,p}}$ multiplies the corresponding measurement
position $(i,p)$ rather than performing an ordinary matrix
multiplication.

% ============================================================
\section{Channel Estimation Problem}
% ============================================================

The matrix measurement model is
\begin{equation}
Y
=
\operatorname{Re}[Z\circ(GS)]
+
|B|
+
N.
\tag{161}
\end{equation}

The objective is to estimate the unknown channel matrix
\begin{equation}
G=[g_{i,k}]
\tag{162}
\end{equation}
from the available measurements.

The known quantities are
\begin{equation}
Y,
\qquad
Z,
\qquad
B,
\qquad
S,
\tag{163}
\end{equation}
while the channel matrix $G$ is unknown.

Let $G^{t}$ denote the current estimate of the channel matrix at
iteration $t$.

Using this estimate, the predicted measurement is
\begin{equation}
\boxed{
\widehat Y(G^{t})
=
\operatorname{Re}[Z\circ(G^{t}S)]
+
|B|
}.
\tag{164}
\end{equation}

The prediction error is defined as
\begin{equation}
\boxed{
E(G^{t})
=
Y-\widehat Y(G^{t})
}.
\tag{165}
\end{equation}

Substituting (164),
\begin{equation}
\boxed{
E(G^{t})
=
Y
-
|B|
-
\operatorname{Re}[Z\circ(G^{t}S)]
}.
\tag{166}
\end{equation}

The dimensions are
\begin{equation}
E(G^{t})
\in
\mathbb{R}^{I\times P}.
\tag{167}
\end{equation}

% ============================================================
\section{Least-Squares Objective}
% ============================================================

The objective is to obtain a channel estimate whose predicted
measurements are as close as possible to the actual measurements.

The squared Frobenius norm is used to measure the total squared
error:
\begin{equation}
\boxed{
\mathcal L(G)
=
\|E(G)\|_F^2
}.
\tag{168}
\end{equation}

\begin{equation}
\boxed{
\mathcal L(G)
=
\left\|
Y-|B|-\operatorname{Re}[Z\circ(GS)]
\right\|_F^2
}.
\tag{169}
\end{equation}

For an $I\times P$ matrix,
\begin{equation}
\|E\|_F^2
=
\sum_{i=1}^{I}
\sum_{p=1}^{P}
|E_{i,p}|^2.
\tag{170}
\end{equation}

Since $E$ is real-valued,
\begin{equation}
\|E\|_F^2
=
\sum_{i=1}^{I}
\sum_{p=1}^{P}
E_{i,p}^2.
\tag{171}
\end{equation}

Therefore, the channel estimate is obtained by solving
\begin{equation}
\boxed{
\widehat G
=
\arg\min_{G}
\mathcal L(G)
}.
\tag{172}
\end{equation}

Equivalently,
\begin{equation}
\boxed{
\widehat G
=
\arg\min_{G}
\left\|
Y-|B|-\operatorname{Re}[Z\circ(GS)]
\right\|_F^2
},
\tag{173}
\end{equation}
which is Eq.~(13) of Xu \emph{et al.}

% ============================================================
\section{Element-by-Element Error Expression}
% ============================================================

From (158) and (166), the $(i,p)$-th error element is
\begin{equation}
E_{i,p}
=
y_{i,p}
-
|b_{i,p}|
-
\operatorname{Re}
\left[
Z_{i,p}
\sum_{r=1}^{K}
g_{i,r}s_{r,p}
\right].
\tag{174}
\end{equation}

Here $r$ is a dummy summation index.

The channel coefficient with respect to which differentiation will
be performed is $g_{i,k}$. The index $k$ is therefore kept fixed
during the differentiation.

The cost function can be written element-wise as
\begin{equation}
\boxed{
\mathcal L
=
\sum_{i=1}^{I}
\sum_{p=1}^{P}
E_{i,p}^2
}.
\tag{175}
\end{equation}

To obtain the gradient with respect to the complex channel matrix,
differentiation is performed with respect to the conjugate variable
$g_{i,k}^*$.

% ============================================================
\section{Wirtinger Differentiation}
% ============================================================

Since $G$ is complex-valued, the loss depends on both $G$ and $G^*$.
Therefore, Wirtinger differentiation is used.

For the selected channel coefficient $g_{i,k}$, only the errors
corresponding to the same receiving-cell index $i$ depend on
$g_{i,k}$. By the chain rule applied to (175),
\begin{equation}
\frac{\partial\mathcal L}{\partial g_{i,k}^*}
=
\sum_{p=1}^{P}
2E_{i,p}
\frac{\partial E_{i,p}}{\partial g_{i,k}^*}.
\tag{176}
\end{equation}

The error element is
\begin{equation}
E_{i,p}
=
y_{i,p}
-
|b_{i,p}|
-
\operatorname{Re}
\left[
Z_{i,p}
\sum_{r=1}^{K}
g_{i,r}s_{r,p}
\right].
\tag{177}
\end{equation}

Using
\begin{equation}
\operatorname{Re}(x)
=
\frac{x+x^*}{2},
\tag{178}
\end{equation}
we obtain
\begin{equation}
E_{i,p}
=
y_{i,p}
-
|b_{i,p}|
-
\frac{1}{2}
\left[
Z_{i,p}
\sum_{r=1}^{K}
g_{i,r}s_{r,p}
+
Z_{i,p}^*
\sum_{r=1}^{K}
g_{i,r}^*s_{r,p}^*
\right].
\tag{179}
\end{equation}

Differentiate with respect to $g_{i,k}^*$:
\begin{equation}
\frac{\partial E_{i,p}}{\partial g_{i,k}^*}
=
-\frac{1}{2}
\frac{\partial}{\partial g_{i,k}^*}
\left[
Z_{i,p}
\sum_{r=1}^{K}
g_{i,r}s_{r,p}
+
Z_{i,p}^*
\sum_{r=1}^{K}
g_{i,r}^*s_{r,p}^*
\right].
\tag{180}
\end{equation}

The first term contains $g_{i,r}$ rather than $g_{i,r}^*$. Under
Wirtinger differentiation,
\begin{equation}
\frac{\partial g_{i,r}}{\partial g_{i,k}^*}
=
0.
\tag{181}
\end{equation}

Therefore,
\begin{equation}
\frac{\partial}{\partial g_{i,k}^*}
\left(
Z_{i,p}
\sum_{r=1}^{K}
g_{i,r}s_{r,p}
\right)
=
0.
\tag{182}
\end{equation}

For the second term,
\begin{equation}
\frac{\partial}{\partial g_{i,k}^*}
\left(
Z_{i,p}^*
\sum_{r=1}^{K}
g_{i,r}^*s_{r,p}^*
\right)
=
Z_{i,p}^*
\sum_{r=1}^{K}
\frac{\partial g_{i,r}^*}{\partial g_{i,k}^*}
s_{r,p}^*.
\tag{183}
\end{equation}

The derivative satisfies
\begin{equation}
\frac{\partial g_{i,r}^*}{\partial g_{i,k}^*}
=
\begin{cases}
1, & r=k,\\
0, & r\neq k.
\end{cases}
\tag{184}
\end{equation}

Consequently,
\begin{equation}
\frac{\partial}{\partial g_{i,k}^*}
\left(
Z_{i,p}^*
\sum_{r=1}^{K}
g_{i,r}^*s_{r,p}^*
\right)
=
Z_{i,p}^*s_{k,p}^*,
\tag{185}
\end{equation}
and
\begin{equation}
\boxed{
\frac{\partial E_{i,p}}{\partial g_{i,k}^*}
=
-\frac{1}{2}
Z_{i,p}^*s_{k,p}^*
}.
\tag{186}
\end{equation}

% ============================================================
\section{Derivative of the Least-Squares Cost}
% ============================================================

Substituting (186) into (176),
\begin{equation}
\begin{aligned}
\frac{\partial\mathcal L}{\partial g_{i,k}^*}
&=
\sum_{p=1}^{P}
2E_{i,p}
\left(
-\frac{1}{2}
Z_{i,p}^*s_{k,p}^*
\right)
=
-\sum_{p=1}^{P}
E_{i,p}
Z_{i,p}^*s_{k,p}^*.
\end{aligned}
\tag{187}
\end{equation}

Hence,
\begin{equation}
\boxed{
\frac{\partial\mathcal L}{\partial g_{i,k}^*}
=
-\sum_{p=1}^{P}
E_{i,p}
Z_{i,p}^*s_{k,p}^*
}.
\tag{188}
\end{equation}

The above operation is performed for every channel coefficient
\begin{equation}
g_{i,k},
\qquad
i=1,\ldots,I,
\qquad
k=1,\ldots,K.
\tag{189}
\end{equation}

Thus, all derivatives are collected into the gradient matrix
\begin{equation}
\boxed{
\frac{\partial\mathcal L}{\partial G^*}
=
\begin{bmatrix}
\dfrac{\partial\mathcal L}{\partial g_{1,1}^*}
&
\dfrac{\partial\mathcal L}{\partial g_{1,2}^*}
&
\cdots
&
\dfrac{\partial\mathcal L}{\partial g_{1,K}^*}
\\[10pt]
\dfrac{\partial\mathcal L}{\partial g_{2,1}^*}
&
\dfrac{\partial\mathcal L}{\partial g_{2,2}^*}
&
\cdots
&
\dfrac{\partial\mathcal L}{\partial g_{2,K}^*}
\\
\vdots&\vdots&\ddots&\vdots
\\
\dfrac{\partial\mathcal L}{\partial g_{I,1}^*}
&
\dfrac{\partial\mathcal L}{\partial g_{I,2}^*}
&
\cdots
&
\dfrac{\partial\mathcal L}{\partial g_{I,K}^*}
\end{bmatrix}
\in
\mathbb{C}^{I\times K}
}.
\tag{190}
\end{equation}

% ============================================================
\section{Conversion of the Gradient to Matrix Form}
% ============================================================

From (188),
\begin{equation}
\frac{\partial\mathcal L}{\partial g_{i,k}^*}
=
-\sum_{p=1}^{P}
E_{i,p}
Z_{i,p}^*s_{k,p}^*.
\tag{191}
\end{equation}

Consider the matrix product
\begin{equation}
(E\circ Z^*)S^H.
\tag{192}
\end{equation}

The dimensions are
\begin{equation}
E\in\mathbb{R}^{I\times P},
\qquad
Z^*\in\mathbb{C}^{I\times P},
\tag{193}
\end{equation}
therefore,
\begin{equation}
E\circ Z^*
\in
\mathbb{C}^{I\times P}.
\tag{194}
\end{equation}

Also,
\begin{equation}
S\in\mathbb{C}^{K\times P},
\qquad
S^H\in\mathbb{C}^{P\times K}.
\tag{195}
\end{equation}

Hence,
\begin{equation}
(E\circ Z^*)S^H:
\qquad
(I\times P)(P\times K)
=
I\times K.
\tag{196}
\end{equation}

Consider the $(i,k)$-th element:
\begin{equation}
[(E\circ Z^*)S^H]_{i,k}
=
\sum_{p=1}^{P}
(E\circ Z^*)_{i,p}
(S^H)_{p,k}.
\tag{197}
\end{equation}

By the definition of the Hadamard product,
\begin{equation}
(E\circ Z^*)_{i,p}
=
E_{i,p}Z_{i,p}^*,
\tag{198}
\end{equation}
and
\begin{equation}
(S^H)_{p,k}
=
S_{k,p}^*
=
s_{k,p}^*.
\tag{199}
\end{equation}

Therefore,
\begin{equation}
[(E\circ Z^*)S^H]_{i,k}
=
\sum_{p=1}^{P}
E_{i,p}Z_{i,p}^*s_{k,p}^*.
\tag{200}
\end{equation}

Comparing (200) with (191),
\begin{equation}
\boxed{
\frac{\partial\mathcal L}{\partial G^*}
=
-(E\circ Z^*)S^H
}.
\tag{201}
\end{equation}

\textbf{Remark (relation to Eq.~(15) of Xu \emph{et al.}).} Xu
\emph{et al.} write
$\nabla_G\mathcal L=-2\big((Y-|B|-\operatorname{Re}(GS\circ Z))\circ
Z^*\big)S^H$, i.e.\ twice (201). The paper uses the convention
$\nabla\mathcal L=2\,\partial\mathcal L/\partial G^*$ (the gradient
with respect to the real and imaginary parts combined as
$\partial/\partial\operatorname{Re}G+j\,\partial/\partial\operatorname{Im}G$),
while (201) is the Wirtinger derivative. Both give the same descent
direction; the step sizes differ by a factor of two.

% ============================================================
\section{Gradient-Descent Channel Update}
% ============================================================

The gradient-descent update rule, Eq.~(14) of Xu \emph{et al.}, is
\begin{equation}
\boxed{
G^{t+1}
=
G^{t}
-
\eta_t
\left.
\frac{\partial\mathcal L}{\partial G^*}
\right|_{G=G^{t}}
},
\tag{202}
\end{equation}
where $\eta_t$ denotes the step size at iteration $t$.

Substituting (201),
\begin{equation}
\begin{aligned}
G^{t+1}
&=
G^{t}
-
\eta_t
\left[
-(E^{t}\circ Z^*)S^H
\right]\\
&=
G^{t}
+
\eta_t
(E^{t}\circ Z^*)S^H.
\end{aligned}
\tag{203}
\end{equation}

The error at iteration $t$ is
\begin{equation}
\boxed{
E^{t}
=
Y
-
|B|
-
\operatorname{Re}[Z\circ(G^{t}S)]
}.
\tag{204}
\end{equation}

Therefore, the complete channel update is
\begin{equation}
\boxed{
G^{t+1}
=
G^{t}
+
\eta_t
\Big(
\big[
Y-|B|-\operatorname{Re}[Z\circ(G^{t}S)]
\big]
\circ Z^*
\Big)
S^H
}.
\tag{205}
\end{equation}

The dimensions of the update are
\begin{equation}
G^{t}\in\mathbb{C}^{I\times K},
\tag{206}
\end{equation}
\begin{equation}
G^{t}S\in\mathbb{C}^{I\times P},
\tag{207}
\end{equation}
\begin{equation}
E^{t}\in\mathbb{R}^{I\times P},
\tag{208}
\end{equation}
\begin{equation}
E^{t}\circ Z^*\in\mathbb{C}^{I\times P},
\tag{209}
\end{equation}
and
\begin{equation}
(E^{t}\circ Z^*)S^H\in\mathbb{C}^{I\times K}.
\tag{210}
\end{equation}

Thus, the gradient has the same dimensions as the channel matrix $G$.

% ============================================================
\section{Gradient-Descent Channel Estimation Algorithm}
% ============================================================

The gradient-descent channel estimation procedure is summarized as
follows.
\begin{equation}
\boxed{
\text{Input: }
Y,\ S,\ Z,\ B,\ G^{0},\ \eta_t,\ \epsilon,\ T_{\max}
}
\tag{211}
\end{equation}

Initialize
\begin{equation}
t=0.
\tag{212}
\end{equation}

At iteration $t$, calculate the predicted measurement:
\begin{equation}
\boxed{
\widehat Y^{t}
=
\operatorname{Re}[Z\circ(G^{t}S)]
+
|B|
}.
\tag{213}
\end{equation}

Calculate the measurement error:
\begin{equation}
\boxed{
E^{t}
=
Y-\widehat Y^{t}
}.
\tag{214}
\end{equation}

Equivalently,
\begin{equation}
\boxed{
E^{t}
=
Y-|B|-\operatorname{Re}[Z\circ(G^{t}S)]
}.
\tag{215}
\end{equation}

Calculate the gradient:
\begin{equation}
\boxed{
D^{t}
=
-(E^{t}\circ Z^*)S^H
}.
\tag{216}
\end{equation}

Update the channel estimate:
\begin{equation}
\boxed{
G^{t+1}
=
G^{t}
-
\eta_t D^{t}
}.
\tag{217}
\end{equation}

Using (216),
\begin{equation}
\boxed{
G^{t+1}
=
G^{t}
+
\eta_t(E^{t}\circ Z^*)S^H
}.
\tag{218}
\end{equation}

The iteration index is then increased according to
\begin{equation}
t\leftarrow t+1.
\tag{219}
\end{equation}

The iterations continue until
\begin{equation}
\boxed{
\left\|
G^{t+1}-G^{t}
\right\|_F^2
<
\epsilon
}
\tag{220}
\end{equation}
or
\begin{equation}
t\geq T_{\max}.
\tag{221}
\end{equation}

The final channel estimate is
\begin{equation}
\boxed{
\widehat G
=
G^{t}
}.
\tag{222}
\end{equation}

% ============================================================
\section{Algorithmic Form}
% ============================================================

The complete gradient-descent procedure is
\begin{equation}
\fbox{%
\begin{minipage}{0.72\linewidth}
\vspace{4pt}
\textbf{Input:} $Y,\ S,\ Z,\ B,\ G^{0},\ \eta_t,\ \epsilon,\ T_{\max}$\\[4pt]
$t=0$\\[4pt]
\textbf{Repeat:}\\[2pt]
\hspace*{1.5em}$\widehat Y^{t}=\operatorname{Re}[Z\circ(G^{t}S)]+|B|$\\[2pt]
\hspace*{1.5em}$E^{t}=Y-\widehat Y^{t}$\\[2pt]
\hspace*{1.5em}$D^{t}=-(E^{t}\circ Z^*)S^H$\\[2pt]
\hspace*{1.5em}$G^{t+1}=G^{t}-\eta_t D^{t}$\\[2pt]
\hspace*{1.5em}$t\leftarrow t+1$\\[4pt]
\textbf{Until:} $\|G^{t+1}-G^{t}\|_F^2<\epsilon$ \ or \ $t\geq T_{\max}$\\[4pt]
\textbf{Output:} $\widehat G=G^{t}$.
\vspace{4pt}
\end{minipage}}
\tag{223}
\end{equation}

% ============================================================
\section{Biased Gerchberg--Saxton Algorithm for the 1D Antenna Array}
% ============================================================

This section follows Algorithm~1 (biased GS) and Section~V
(extension to channel estimation) of Cui \emph{et al.},
``Towards Atomic MIMO Receivers.''

\subsection{1D Channel Estimation Model}

Consider an array consisting of $I$ atomic vapor cells and $K$ users.

For the $i$-th vapor cell, define the channel vector as
\begin{equation}
\boxed{
\mathbf g_i
=
\begin{bmatrix}
g_{i,1}\\
g_{i,2}\\
\vdots\\
g_{i,K}
\end{bmatrix}
\in
\mathbb{C}^{K\times1}
}.
\tag{224}
\end{equation}

The known pilot matrix is
\begin{equation}
\boxed{
S=
\begin{bmatrix}
s_{1,1}&s_{1,2}&\cdots&s_{1,P}\\
s_{2,1}&s_{2,2}&\cdots&s_{2,P}\\
\vdots&\vdots&\ddots&\vdots\\
s_{K,1}&s_{K,2}&\cdots&s_{K,P}
\end{bmatrix}
\in
\mathbb{C}^{K\times P}
}.
\tag{225}
\end{equation}

The reference vector and noise vector are
\begin{equation}
\mathbf b_i\in\mathbb{C}^{P\times1},
\qquad
\mathbf w_i\in\mathbb{C}^{P\times1}.
\tag{226}
\end{equation}

The measured magnitude vector, i.e.\ the $i$-th row of the exact
model (73) written as a column, is
\begin{equation}
\boxed{
\mathbf y_i
=
\left|
S^T\mathbf g_i+\mathbf b_i+\mathbf w_i
\right|
\in
\mathbb{R}^{P\times1}
}.
\tag{227}
\end{equation}

The $p$-th element of $S^T\mathbf g_i$ is
\begin{equation}
\left[S^T\mathbf g_i\right]_p
=
\sum_{k=1}^{K}
s_{k,p}g_{i,k}.
\tag{228}
\end{equation}

The difficulty is that only the magnitude of the received complex
signal is measured.

Define the unknown complex received vector
\begin{equation}
\boxed{
\mathbf x_i
=
S^T\mathbf g_i+\mathbf b_i+\mathbf w_i
}.
\tag{229}
\end{equation}

Then,
\begin{equation}
\mathbf y_i=|\mathbf x_i|.
\tag{230}
\end{equation}

If the phase of $\mathbf x_i$ were known, the complex received signal
could be reconstructed as
\begin{equation}
\mathbf x_i
=
\mathbf y_i\circ e^{j\boldsymbol\theta_i},
\tag{231}
\end{equation}
where $\boldsymbol\theta_i$ denotes the unknown phase vector.

The GS procedure therefore alternates between phase estimation,
complex-signal reconstruction, and channel estimation. Unlike the GD
method of the previous sections, GS works with the exact magnitude
model (227) and does not use the strong-reference linearization.

% ============================================================
\section{Spectral Initialization}
% ============================================================

The iterative GS algorithm depends on its initial channel estimate.
A spectral initialization is therefore used to obtain an initial
estimate of the channel.

The biased GS initialization constructs an augmented observation
model, forms a weighted covariance matrix, extracts its principal
eigenvector, estimates the corresponding magnitude, and removes the
arbitrary global phase.

\subsection{Conversion to the Standard Spectral Form}

The channel measurement model is
\begin{equation}
\mathbf y_i
=
\left|
S^T\mathbf g_i+\mathbf b_i+\mathbf w_i
\right|.
\tag{232}
\end{equation}

The standard biased phase-retrieval model of Cui \emph{et al.}\ is
written as
\begin{equation}
\mathbf z
=
\left|
A^H\mathbf g+\mathbf b+\mathbf w
\right|.
\tag{233}
\end{equation}

Comparing (232) and (233),
\begin{equation}
\boxed{
A^H=S^T
}.
\tag{234}
\end{equation}

Taking the Hermitian relation,
\begin{equation}
A=S^*,
\tag{235}
\end{equation}
since $(S^*)^H=S^T$.

Therefore, the channel model can be written as
\begin{equation}
\boxed{
\mathbf y_i
=
\left|
A^H\mathbf g_i+\mathbf b_i+\mathbf w_i
\right|,
\qquad
A=S^*
}.
\tag{236}
\end{equation}

This mapping permits the standard spectral initialization procedure
to be applied without changing the definition of the channel vector
$\mathbf g_i$.

\subsection{Augmented Channel Vector}

The received signal contains an unknown and a known component,
\begin{equation}
A^H\mathbf g_i+\mathbf b_i.
\tag{237}
\end{equation}

The unknown channel and the known reference are combined into one
augmented vector:
\begin{equation}
\boxed{
\bar{\mathbf g}_i
=
\begin{bmatrix}
\mathbf g_i\\
1
\end{bmatrix}
\in
\mathbb{C}^{(K+1)\times1}
}.
\tag{238}
\end{equation}

The last element is fixed to one so that
\begin{equation}
\begin{bmatrix}
A^H&\mathbf b_i
\end{bmatrix}
\begin{bmatrix}
\mathbf g_i\\
1
\end{bmatrix}
=
A^H\mathbf g_i+\mathbf b_i.
\tag{239}
\end{equation}

Define the augmented observation matrix as
\begin{equation}
\boxed{
\bar A_i
=
\begin{bmatrix}
A^H&\mathbf b_i
\end{bmatrix}
=
\begin{bmatrix}
S^T&\mathbf b_i
\end{bmatrix}
}.
\tag{240}
\end{equation}

\textbf{Remark.} Cui \emph{et al.}\ define
$\bar A=[A^H,\mathbf b]^H\in\mathbb C^{(K+1)\times N}$. The matrix
$\bar A_i$ in (240) is therefore their $\bar A^H$; the columns
$\bar{\mathbf a}_{i,p}$ of $\bar A_i^H$ defined below coincide with
their $\bar{\mathbf a}_n$.

The dimensions are
\begin{equation}
S^T\in\mathbb{C}^{P\times K},
\qquad
\mathbf b_i\in\mathbb{C}^{P\times1},
\tag{241}
\end{equation}
and hence
\begin{equation}
\boxed{
\bar A_i
\in
\mathbb{C}^{P\times(K+1)}
}.
\tag{242}
\end{equation}

The measurement model becomes
\begin{equation}
\boxed{
\mathbf y_i
=
\left|
\bar A_i\bar{\mathbf g}_i+\mathbf w_i
\right|
}.
\tag{243}
\end{equation}

\subsection{Construction of the Augmented Observation Vectors}

The Hermitian transpose of the augmented matrix is
\begin{equation}
\bar A_i^H
\in
\mathbb{C}^{(K+1)\times P}.
\tag{244}
\end{equation}

Its $p$-th column is
\begin{equation}
\boxed{
\bar{\mathbf a}_{i,p}
=
\begin{bmatrix}
s_{1,p}^*\\
s_{2,p}^*\\
\vdots\\
s_{K,p}^*\\
b_{i,p}^*
\end{bmatrix}
\in
\mathbb{C}^{(K+1)\times1}
}.
\tag{245}
\end{equation}

Since
\begin{equation}
\bar A_i^H
=
\begin{bmatrix}
\bar{\mathbf a}_{i,1}&
\bar{\mathbf a}_{i,2}&
\cdots&
\bar{\mathbf a}_{i,P}
\end{bmatrix},
\tag{246}
\end{equation}
the corresponding inner product is
\begin{equation}
\bar{\mathbf a}_{i,p}^H\bar{\mathbf g}_i
=
\sum_{k=1}^{K}
s_{k,p}g_{i,k}
+
b_{i,p}.
\tag{247}
\end{equation}

Thus,
\begin{equation}
\boxed{
\bar{\mathbf a}_{i,p}^H\bar{\mathbf g}_i
=
\left[S^T\mathbf g_i\right]_p
+
b_{i,p}
}.
\tag{248}
\end{equation}

% ============================================================
\section{Spectral Matrix Construction}
% ============================================================

\subsection{Weighted Covariance Matrix}

Following Step~2 of Algorithm~1 in Cui \emph{et al.}, the spectral
method constructs the weighted covariance matrix
\begin{equation}
\boxed{
M_i
=
\sum_{p=1}^{P}
y_{i,p}
\bar{\mathbf a}_{i,p}
\bar{\mathbf a}_{i,p}^H
}.
\tag{249}
\end{equation}

The dimensions are
\begin{equation}
\bar{\mathbf a}_{i,p}
\in
\mathbb{C}^{(K+1)\times1},
\tag{250}
\end{equation}
and therefore
\begin{equation}
\bar{\mathbf a}_{i,p}
\bar{\mathbf a}_{i,p}^H
\in
\mathbb{C}^{(K+1)\times(K+1)}.
\tag{251}
\end{equation}

Consequently,
\begin{equation}
\boxed{
M_i
\in
\mathbb{C}^{(K+1)\times(K+1)}
}.
\tag{252}
\end{equation}

The matrix $M_i$ is Hermitian, since each term
$y_{i,p}\bar{\mathbf a}_{i,p}\bar{\mathbf a}_{i,p}^H$ is Hermitian and
$y_{i,p}\in\mathbb R$. It contains the magnitude measurements
$y_{i,p}$ as weights for the corresponding observation-vector outer
products.

\subsection{Principal Eigenvector}

Perform the eigendecomposition
\begin{equation}
\boxed{
M_i\mathbf v_{i,r}
=
\lambda_{i,r}\mathbf v_{i,r}
}.
\tag{253}
\end{equation}

Select the eigenvector associated with the largest eigenvalue:
\begin{equation}
\boxed{
\mathbf v_i
=
\mathbf v_{i,r_{\max}}
}.
\tag{254}
\end{equation}

Equivalently,
\begin{equation}
M_i\mathbf v_i
=
\lambda_{\max}(M_i)\mathbf v_i.
\tag{255}
\end{equation}

The principal eigenvector satisfies
\begin{equation}
\boxed{
\mathbf v_i
\in
\mathbb{C}^{(K+1)\times1},
\qquad
\|\mathbf v_i\|_2=1
}.
\tag{256}
\end{equation}

The spectral method uses the principal eigenvector as an estimate of
the direction of the augmented channel vector. Thus,
\begin{equation}
\mathbf v_i
\approx
e^{j\phi_i}
\frac{\bar{\mathbf g}_i}{\|\bar{\mathbf g}_i\|_2},
\tag{257}
\end{equation}
where $\phi_i$ represents an arbitrary phase.

The eigenvector provides the direction of the augmented channel, but
its magnitude is not yet determined.

% ============================================================
\section{Spectral Magnitude Estimation}
% ============================================================

Assume the initial augmented channel vector has the form
\begin{equation}
\boxed{
\bar{\mathbf g}_i^{(0)}
=
r_i\mathbf v_i
},
\tag{258}
\end{equation}
where $r_i$ is a real nonnegative scalar.

The predicted magnitude corresponding to this estimate is
\begin{equation}
\left|
\bar A_i\bar{\mathbf g}_i^{(0)}
\right|
=
\left|
r_i\bar A_i\mathbf v_i
\right|.
\tag{259}
\end{equation}

Since $r_i\geq0$,
\begin{equation}
\left|
r_i\bar A_i\mathbf v_i
\right|
=
r_i
\left|
\bar A_i\mathbf v_i
\right|.
\tag{260}
\end{equation}

The magnitude should approximate the measured vector $\mathbf y_i$.
Therefore,
\begin{equation}
\boxed{
r_i
=
\arg\min_{r_i}
\left\|
r_i\left|\bar A_i\mathbf v_i\right|
-
\mathbf y_i
\right\|_2^2
}.
\tag{261}
\end{equation}

Define
\begin{equation}
\mathbf q_i
=
\left|
\bar A_i\mathbf v_i
\right|
\in\mathbb R^{P\times1}.
\tag{262}
\end{equation}

Then the scalar least-squares problem becomes
\begin{equation}
r_i
=
\arg\min_{r_i}
\left\|
r_i\mathbf q_i-\mathbf y_i
\right\|_2^2.
\tag{263}
\end{equation}

Expanding the squared norm, with $\mathbf q_i$ and $\mathbf y_i$ real,
\begin{equation}
\left\|
r_i\mathbf q_i-\mathbf y_i
\right\|_2^2
=
r_i^2\mathbf q_i^T\mathbf q_i
-
2r_i\mathbf q_i^T\mathbf y_i
+
\mathbf y_i^T\mathbf y_i.
\tag{264}
\end{equation}

Since $r_i$ is real,
\begin{equation}
\frac{\partial}{\partial r_i}
\left\|
r_i\mathbf q_i-\mathbf y_i
\right\|_2^2
=
2r_i\|\mathbf q_i\|_2^2
-
2\mathbf q_i^T\mathbf y_i.
\tag{265}
\end{equation}

Setting the derivative equal to zero,
\begin{equation}
2r_i\|\mathbf q_i\|_2^2
-
2\mathbf q_i^T\mathbf y_i
=
0.
\tag{266}
\end{equation}

Therefore,
\begin{equation}
r_i\|\mathbf q_i\|_2^2
=
\mathbf q_i^T\mathbf y_i,
\tag{267}
\end{equation}
and
\begin{equation}
\boxed{
r_i
=
\frac{\mathbf q_i^T\mathbf y_i}{\|\mathbf q_i\|_2^2}
}.
\tag{268}
\end{equation}

Using (262),
\begin{equation}
\boxed{
r_i
=
\frac{
\left|\bar A_i\mathbf v_i\right|^T\mathbf y_i
}{
\left\|\bar A_i\mathbf v_i\right\|_2^2
}
},
\tag{269}
\end{equation}
which is Step~3 of Algorithm~1 in Cui \emph{et al.}, since
$|\mathbf v^H\bar A|=|\bar A^H\mathbf v|^T$ in their notation.

The initial augmented channel estimate is therefore
\begin{equation}
\boxed{
\bar{\mathbf g}_i^{(0)}
=
r_i\mathbf v_i
}.
\tag{270}
\end{equation}

% ============================================================
\section{Global Phase Removal}
% ============================================================

Magnitude measurements cannot distinguish between
\begin{equation}
\bar{\mathbf g}_i
\tag{271}
\end{equation}
and
\begin{equation}
e^{j\phi_i}\bar{\mathbf g}_i,
\tag{272}
\end{equation}
because
\begin{equation}
\left|
\bar A_ie^{j\phi_i}\bar{\mathbf g}_i
\right|
=
\left|
e^{j\phi_i}\bar A_i\bar{\mathbf g}_i
\right|
=
\left|
\bar A_i\bar{\mathbf g}_i
\right|.
\tag{273}
\end{equation}

However, the augmented vector is defined as
\begin{equation}
\bar{\mathbf g}_i
=
\begin{bmatrix}
\mathbf g_i\\
1
\end{bmatrix},
\tag{274}
\end{equation}
so the final element must have zero phase.

Let
\begin{equation}
\boxed{
\phi_i
=
\angle
\left(
[\bar{\mathbf g}_i^{(0)}]_{K+1}
\right)
}.
\tag{275}
\end{equation}

Multiply the entire augmented vector by
\begin{equation}
e^{-j\phi_i}.
\tag{276}
\end{equation}

Thus,
\begin{equation}
\boxed{
\widetilde{\mathbf g}_i^{(0)}
=
e^{-j\phi_i}\bar{\mathbf g}_i^{(0)}
}.
\tag{277}
\end{equation}

The last element becomes
\begin{equation}
[\widetilde{\mathbf g}_i^{(0)}]_{K+1}
=
e^{-j\phi_i}
[\bar{\mathbf g}_i^{(0)}]_{K+1}.
\tag{278}
\end{equation}

Using (275),
\begin{equation}
[\bar{\mathbf g}_i^{(0)}]_{K+1}
=
\left|
[\bar{\mathbf g}_i^{(0)}]_{K+1}
\right|
e^{j\phi_i},
\tag{279}
\end{equation}
and therefore,
\begin{equation}
\begin{aligned}
[\widetilde{\mathbf g}_i^{(0)}]_{K+1}
&=
e^{-j\phi_i}
\left|
[\bar{\mathbf g}_i^{(0)}]_{K+1}
\right|
e^{j\phi_i}
=
\left|
[\bar{\mathbf g}_i^{(0)}]_{K+1}
\right|.
\end{aligned}
\tag{280}
\end{equation}

Hence, the final element has zero phase.

The channel estimate for the $i$-th antenna is obtained by retaining
the first $K$ elements (Step~4 of Algorithm~1 in Cui \emph{et al.}):
\begin{equation}
\boxed{
\mathbf g_i^{(0)}
=
\widetilde{\mathbf g}_i^{(0)}(1:K)
}.
\tag{281}
\end{equation}

\subsection{Construction of the Initial Channel Matrix}

The above spectral initialization is performed independently for all
$I$ vapor cells. For
\begin{equation}
i=1,2,\ldots,I,
\tag{282}
\end{equation}
obtain
\begin{equation}
\mathbf g_i^{(0)}
\in
\mathbb{C}^{K\times1}.
\tag{283}
\end{equation}

The initial channel matrix is constructed by placing the channel
vectors as rows:
\begin{equation}
\boxed{
G^{(0)}
=
\begin{bmatrix}
(\mathbf g_1^{(0)})^T\\
(\mathbf g_2^{(0)})^T\\
\vdots\\
(\mathbf g_I^{(0)})^T
\end{bmatrix}
\in
\mathbb{C}^{I\times K}
}.
\tag{284}
\end{equation}

This matrix provides the initial channel estimate for the subsequent
GS iterations.

% ============================================================
\section{Gerchberg--Saxton Iteration}
% ============================================================

\subsection{Step 1: Prediction of the Complex Received Signal}

Suppose that the channel estimate from the previous iteration is
\begin{equation}
G^{(t-1)}
\in
\mathbb{C}^{I\times K}.
\tag{285}
\end{equation}

For the $i$-th antenna,
\begin{equation}
\mathbf g_i^{(t-1)}
\in
\mathbb{C}^{K\times1}.
\tag{286}
\end{equation}

The predicted complex received signal is
\begin{equation}
\boxed{
\mathbf x_i^{(t)}
=
S^T\mathbf g_i^{(t-1)}
+
\mathbf b_i
}.
\tag{287}
\end{equation}

The noise realization is not added to the predicted signal because
the actual noise realization is unknown. The measured magnitude
already contains the noise contribution.

The dimensions are
\begin{equation}
S^T\in\mathbb{C}^{P\times K},
\qquad
\mathbf g_i^{(t-1)}\in\mathbb{C}^{K\times1},
\tag{288}
\end{equation}
hence
\begin{equation}
\mathbf x_i^{(t)}
\in
\mathbb{C}^{P\times1}.
\tag{289}
\end{equation}

\subsection{Step 2: Phase Estimation}

The phase of the predicted complex received signal (Step~6 of
Algorithm~1 in Cui \emph{et al.}) is
\begin{equation}
\boxed{
\boldsymbol\theta_i^{(t)}
=
\angle
\left(
S^T\mathbf g_i^{(t-1)}
+
\mathbf b_i
\right)
}.
\tag{290}
\end{equation}

Collecting the phase values gives
\begin{equation}
\boldsymbol\theta_i^{(t)}
=
\begin{bmatrix}
\theta_{i,1}^{(t)}\\
\theta_{i,2}^{(t)}\\
\vdots\\
\theta_{i,P}^{(t)}
\end{bmatrix}
\in
\mathbb{R}^{P\times1}.
\tag{291}
\end{equation}

The corresponding unit-magnitude phase vector is
\begin{equation}
e^{j\boldsymbol\theta_i^{(t)}}.
\tag{292}
\end{equation}

Each element has unit magnitude:
\begin{equation}
\left|
e^{j\theta_{i,p}^{(t)}}
\right|
=
1.
\tag{293}
\end{equation}

\subsection{Step 3: Replace the Predicted Magnitude with the Measured Magnitude}

The measured vector is
\begin{equation}
\mathbf y_i
=
|\mathbf x_i|.
\tag{294}
\end{equation}

The GS procedure retains the measured magnitude and replaces only the
unknown phase by the current phase estimate. Define
\begin{equation}
\boxed{
\mathbf r_i^{(t)}
=
\mathbf y_i\circ e^{j\boldsymbol\theta_i^{(t)}}
}.
\tag{295}
\end{equation}

Element-wise,
\begin{equation}
\boxed{
r_{i,p}^{(t)}
=
y_{i,p}e^{j\theta_{i,p}^{(t)}}
}.
\tag{296}
\end{equation}

Therefore,
\begin{equation}
\left|
r_{i,p}^{(t)}
\right|
=
y_{i,p}
\left|
e^{j\theta_{i,p}^{(t)}}
\right|
=
y_{i,p}.
\tag{297}
\end{equation}

Thus,
\begin{equation}
\boxed{
|\mathbf r_i^{(t)}|
=
\mathbf y_i
}.
\tag{298}
\end{equation}

This operation replaces the unknown complex received signal with a
complex vector having the measured magnitude and the currently
estimated phase.

% ============================================================
\section{Reference Removal and Least-Squares Channel Update}
% ============================================================

\subsection{Step 4: Remove the Known Reference}

The model for the reconstructed received signal is
\begin{equation}
\mathbf r_i^{(t)}
\approx
S^T\mathbf g_i+\mathbf b_i.
\tag{299}
\end{equation}

Therefore,
\begin{equation}
\boxed{
\mathbf r_i^{(t)}-\mathbf b_i
\approx
S^T\mathbf g_i
}.
\tag{300}
\end{equation}

After removing the known reference, the channel estimation becomes a
standard linear least-squares problem.

\subsection{Step 5: Least-Squares Channel Estimation}

The channel vector is obtained from
\begin{equation}
\boxed{
\mathbf g_i^{(t)}
=
\arg\min_{\mathbf g}
\left\|
\mathbf r_i^{(t)}-\mathbf b_i-S^T\mathbf g
\right\|_2^2
}.
\tag{301}
\end{equation}

Using $A=S^*$ and $A^H=S^T$,
\begin{equation}
\mathbf g_i^{(t)}
=
\arg\min_{\mathbf g}
\left\|
(\mathbf r_i^{(t)}-\mathbf b_i)-A^H\mathbf g
\right\|_2^2.
\tag{302}
\end{equation}

The standard least-squares solution (Step~7 of Algorithm~1 in Cui
\emph{et al.}) is
\begin{equation}
\boxed{
\mathbf g_i^{(t)}
=
(AA^H)^{-1}A
\left(
\mathbf r_i^{(t)}-\mathbf b_i
\right)
}.
\tag{303}
\end{equation}

Since
\begin{equation}
AA^H
=
S^*S^T,
\tag{304}
\end{equation}
we have
\begin{equation}
\boxed{
\mathbf g_i^{(t)}
=
(S^*S^T)^{-1}S^*
\left(
\mathbf r_i^{(t)}-\mathbf b_i
\right)
}.
\tag{305}
\end{equation}

Using (295),
\begin{equation}
\boxed{
\mathbf g_i^{(t)}
=
(S^*S^T)^{-1}S^*
\left(
\mathbf y_i\circ e^{j\boldsymbol\theta_i^{(t)}}
-
\mathbf b_i
\right)
}.
\tag{306}
\end{equation}

The inverse exists when $S$ has full row rank, which requires
$P\geq K$.

% ============================================================
\section{Matrix Form of the GS Iteration}
% ============================================================

\subsection{Simultaneous Update for All Antennas}

Instead of processing each antenna separately, define
\begin{equation}
\boxed{
G^{(t-1)}
\in
\mathbb{C}^{I\times K}
}.
\tag{307}
\end{equation}

The predicted complex received matrix is
\begin{equation}
\boxed{
X^{(t)}
=
G^{(t-1)}S+B
}.
\tag{308}
\end{equation}

The dimensions are
\begin{equation}
G^{(t-1)}S:
\qquad
(I\times K)(K\times P)
=
I\times P,
\tag{309}
\end{equation}
and
\begin{equation}
B\in\mathbb{C}^{I\times P}.
\tag{310}
\end{equation}

Hence,
\begin{equation}
\boxed{
X^{(t)}
\in
\mathbb{C}^{I\times P}
}.
\tag{311}
\end{equation}

The phase matrix is
\begin{equation}
\boxed{
\Theta^{(t)}
=
\angle X^{(t)}
}.
\tag{312}
\end{equation}

The reconstructed complex measurement matrix is
\begin{equation}
\boxed{
R^{(t)}
=
Y\circ e^{j\Theta^{(t)}}
}.
\tag{313}
\end{equation}

Since $Y$ contains the measured magnitudes,
\begin{equation}
|R^{(t)}|
=
Y.
\tag{314}
\end{equation}

The known reference matrix is then removed:
\begin{equation}
\boxed{
R^{(t)}-B
}.
\tag{315}
\end{equation}

The channel update is obtained by solving
\begin{equation}
\boxed{
G^{(t)}
=
\arg\min_{G}
\left\|
R^{(t)}-B-GS
\right\|_F^2
}.
\tag{316}
\end{equation}

The least-squares solution is
\begin{equation}
\boxed{
G^{(t)}
=
(R^{(t)}-B)S^H(SS^H)^{-1}
}.
\tag{317}
\end{equation}

The dimensions are
\begin{equation}
R^{(t)}-B\in\mathbb{C}^{I\times P},
\tag{318}
\end{equation}
\begin{equation}
S^H\in\mathbb{C}^{P\times K},
\tag{319}
\end{equation}
and
\begin{equation}
SS^H\in\mathbb{C}^{K\times K}.
\tag{320}
\end{equation}

Therefore,
\begin{equation}
(R^{(t)}-B)S^H
\in
\mathbb{C}^{I\times K},
\tag{321}
\end{equation}
and
\begin{equation}
\boxed{
G^{(t)}
\in
\mathbb{C}^{I\times K}
}.
\tag{322}
\end{equation}

\subsection{Equivalence with the Per-Antenna LS Solution}

For the $i$-th row of the matrix model,
\begin{equation}
[R^{(t)}-B]_{i,:}
=
(\mathbf r_i^{(t)}-\mathbf b_i)^T.
\tag{323}
\end{equation}

The corresponding row of $G^{(t)}$ is
\begin{equation}
[G^{(t)}]_{i,:}
=
(\mathbf g_i^{(t)})^T.
\tag{324}
\end{equation}

From (317),
\begin{equation}
[G^{(t)}]_{i,:}
=
[R^{(t)}-B]_{i,:}
S^H(SS^H)^{-1}.
\tag{325}
\end{equation}

Taking the transpose and using $(SS^H)^T=S^*S^T$ and
$(S^H)^T=S^*$ gives the corresponding channel-vector representation,
\begin{equation}
\mathbf g_i^{(t)}
=
(S^*S^T)^{-1}S^*
(\mathbf r_i^{(t)}-\mathbf b_i),
\tag{326}
\end{equation}
which is (305).

Thus, the matrix update and the per-antenna least-squares solution
represent the same channel-estimation operation under different
orientations of the channel vectors.

% ============================================================
\section{Convergence Criterion}
% ============================================================

\subsection{Difference Between Consecutive Channel Estimates}

After obtaining $G^{(t)}$, compare it with the channel estimate from
the previous iteration:
\begin{equation}
\boxed{
\Delta_t
=
\left\|
G^{(t)}-G^{(t-1)}
\right\|_F^2
}.
\tag{327}
\end{equation}

If
\begin{equation}
\boxed{
\Delta_t<\epsilon
},
\tag{328}
\end{equation}
where $\epsilon$ is a prescribed convergence tolerance, the iteration
is terminated. Otherwise, the GS iterations continue.

A maximum iteration number may also be imposed as a stopping
condition:
\begin{equation}
\boxed{
t\geq T_{\max}
}.
\tag{329}
\end{equation}

Therefore, the algorithm terminates when either
\begin{equation}
\Delta_t<\epsilon
\tag{330}
\end{equation}
or
\begin{equation}
t\geq T_{\max}.
\tag{331}
\end{equation}

\textbf{Remark.} Cui \emph{et al.}\ use a fixed number of iterations
($t_0=50$); the tolerance test (328) is an addition that leaves the
fixed point unchanged.

% ============================================================
\section{Final Biased GS Algorithm for the 1D Array}
% ============================================================

The complete biased GS channel-estimation algorithm consists of
spectral initialization followed by iterative phase estimation and
least-squares channel updates.

\subsection{Algorithm Input}

The required inputs are
\begin{equation}
\boxed{
Y,\quad S,\quad B,\quad T_{\max},\quad \epsilon
}.
\tag{332}
\end{equation}

For the per-antenna spectral initialization,
\begin{equation}
\mathbf y_i\in\mathbb{R}^{P\times1},
\tag{333}
\end{equation}
\begin{equation}
S\in\mathbb{C}^{K\times P},
\tag{334}
\end{equation}
and
\begin{equation}
\mathbf b_i\in\mathbb{C}^{P\times1}.
\tag{335}
\end{equation}

\subsection{Spectral Initialization Procedure}

For
\begin{equation}
i=1,2,\ldots,I,
\tag{336}
\end{equation}
construct
\begin{equation}
\boxed{
\bar A_i
=
\begin{bmatrix}
S^T&\mathbf b_i
\end{bmatrix}
}.
\tag{337}
\end{equation}

Let $\bar{\mathbf a}_{i,p}$ denote the $p$-th column of $\bar A_i^H$.
Construct
\begin{equation}
\boxed{
M_i
=
\sum_{p=1}^{P}
y_{i,p}
\bar{\mathbf a}_{i,p}
\bar{\mathbf a}_{i,p}^H
}.
\tag{338}
\end{equation}

Obtain the principal eigenvector
\begin{equation}
\boxed{
\mathbf v_i
=
\operatorname{PrincipalEigenvector}(M_i)
}.
\tag{339}
\end{equation}

Estimate the magnitude scaling factor:
\begin{equation}
\boxed{
r_i
=
\frac{
|\bar A_i\mathbf v_i|^T\mathbf y_i
}{
\|\bar A_i\mathbf v_i\|_2^2
}
}.
\tag{340}
\end{equation}

Construct the initial augmented channel vector:
\begin{equation}
\boxed{
\bar{\mathbf g}_i^{(0)}
=
r_i\mathbf v_i
}.
\tag{341}
\end{equation}

Remove its arbitrary global phase:
\begin{equation}
\boxed{
\widetilde{\mathbf g}_i^{(0)}
=
e^{-j\angle([\bar{\mathbf g}_i^{(0)}]_{K+1})}
\bar{\mathbf g}_i^{(0)}
}.
\tag{342}
\end{equation}

Retain the first $K$ elements:
\begin{equation}
\boxed{
\mathbf g_i^{(0)}
=
\widetilde{\mathbf g}_i^{(0)}(1:K)
}.
\tag{343}
\end{equation}

After processing all antennas, construct
\begin{equation}
\boxed{
G^{(0)}
=
\begin{bmatrix}
(\mathbf g_1^{(0)})^T\\
(\mathbf g_2^{(0)})^T\\
\vdots\\
(\mathbf g_I^{(0)})^T
\end{bmatrix}
}.
\tag{344}
\end{equation}

\subsection{GS Iterations}

For
\begin{equation}
t=1,2,\ldots,T_{\max},
\tag{345}
\end{equation}
calculate the predicted complex received matrix:
\begin{equation}
\boxed{
X^{(t)}
=
G^{(t-1)}S+B
}.
\tag{346}
\end{equation}

Calculate its phase:
\begin{equation}
\boxed{
\Theta^{(t)}
=
\angle X^{(t)}
}.
\tag{347}
\end{equation}

Replace the predicted magnitude by the measured magnitude:
\begin{equation}
\boxed{
R^{(t)}
=
Y\circ e^{j\Theta^{(t)}}
}.
\tag{348}
\end{equation}

Remove the reference:
\begin{equation}
\boxed{
R^{(t)}-B
}.
\tag{349}
\end{equation}

Update the channel using least squares:
\begin{equation}
\boxed{
G^{(t)}
=
(R^{(t)}-B)S^H(SS^H)^{-1}
}.
\tag{350}
\end{equation}

Evaluate the convergence measure:
\begin{equation}
\boxed{
\Delta_t
=
\left\|
G^{(t)}-G^{(t-1)}
\right\|_F^2
}.
\tag{351}
\end{equation}

If
\begin{equation}
\boxed{
\Delta_t\leq\epsilon,
}
\tag{352}
\end{equation}
terminate the iterations. Otherwise, continue until
\begin{equation}
t=T_{\max}.
\tag{353}
\end{equation}

The final channel estimate is
\begin{equation}
\boxed{
\widehat G
=
G^{(t)}
}.
\tag{354}
\end{equation}

% ============================================================
\section{Complete Algorithm}
% ============================================================

The complete biased GS algorithm for the 1D atomic receiver array can
therefore be summarized as
\begin{equation}
\fbox{%
\begin{minipage}{0.72\linewidth}
\vspace{4pt}
\textbf{Input:} $Y,\ S,\ B,\ T_{\max},\ \epsilon$\\[6pt]
\textbf{Spectral Initialization:}\\[2pt]
\hspace*{1.5em}For $i=1,2,\ldots,I$:\\[2pt]
\hspace*{3em}$\bar A_i=[S^T\ \ \mathbf b_i]$\\[2pt]
\hspace*{3em}$\bar{\mathbf a}_{i,p}=p$-th column of $\bar A_i^H$\\[2pt]
\hspace*{3em}$M_i=\displaystyle\sum_{p=1}^{P}y_{i,p}\bar{\mathbf a}_{i,p}\bar{\mathbf a}_{i,p}^H$\\[2pt]
\hspace*{3em}$\mathbf v_i=$ principal eigenvector of $M_i$\\[2pt]
\hspace*{3em}$r_i=\dfrac{|\bar A_i\mathbf v_i|^T\mathbf y_i}{\|\bar A_i\mathbf v_i\|_2^2}$\\[4pt]
\hspace*{3em}$\bar{\mathbf g}_i^{(0)}=r_i\mathbf v_i$\\[2pt]
\hspace*{3em}$\widetilde{\mathbf g}_i^{(0)}=e^{-j\angle([\bar{\mathbf g}_i^{(0)}]_{K+1})}\bar{\mathbf g}_i^{(0)}$\\[2pt]
\hspace*{3em}$\mathbf g_i^{(0)}=\widetilde{\mathbf g}_i^{(0)}(1:K)$\\[2pt]
\hspace*{1.5em}End\\[2pt]
\hspace*{1.5em}$G^{(0)}=\big[(\mathbf g_1^{(0)})^T;\ \ldots;\ (\mathbf g_I^{(0)})^T\big]$\\[6pt]
\textbf{GS Iterations:}\\[2pt]
\hspace*{1.5em}For $t=1,2,\ldots,T_{\max}$:\\[2pt]
\hspace*{3em}$X^{(t)}=G^{(t-1)}S+B$\\[2pt]
\hspace*{3em}$\Theta^{(t)}=\angle X^{(t)}$\\[2pt]
\hspace*{3em}$R^{(t)}=Y\circ e^{j\Theta^{(t)}}$\\[2pt]
\hspace*{3em}$G^{(t)}=(R^{(t)}-B)S^H(SS^H)^{-1}$\\[2pt]
\hspace*{3em}$\Delta_t=\|G^{(t)}-G^{(t-1)}\|_F^2$\\[2pt]
\hspace*{3em}If $\Delta_t\leq\epsilon$: stop\\[2pt]
\hspace*{1.5em}End\\[6pt]
\textbf{Output:} $\widehat G=G^{(t)}$.
\vspace{4pt}
\end{minipage}}
\tag{355}
\end{equation}

\end{document}
